%% LyX 2.1.3 created this file.  For more info, see http://www.lyx.org/.
%% Do not edit unless you really know what you are doing.
\ifdefined\directlua
  \documentclass[english,french,latin,german,twoside,a4paper]{ltjsarticle}
\else
  \documentclass[english,french,latin,german,twoside,a4paper]{scrartcl}
\fi
\usepackage{version}
\usepackage{tabularx}
\usepackage{tcolorbox}
\usepackage{amsfonts}
\usepackage{todonotes}
\usepackage{metalogo}

\includeversion{frakturfonttypesetting}
\excludeversion{romanfonttypesetting}

\newif\ifFraktur
\Frakturtrue
%\Frakturfalse

\input{preamble.tex}

\setstretch{1.375}

\hypersetup{pdftitle={Sprachen Sie Attich? 日本語訳第2部},
  pdfauthor={J. Sengbusch}}


\usepackage{polyglossia}
\setotherlanguage[spelling=old,babelshorthands=true]{german}
\setotherlanguage{english}
\setotherlanguage{french}
\setotherlanguage[variant=ancient]{greek}
\setotherlanguage{latin}
\setmainlanguage{japanese}

\newcommand{\trDE}[1]{\texorpdfstring{\begin{german}#1\end{german}}{#1}}
\newcommand{\trLA}[1]{\texorpdfstring{}{}}
\newcommand{\trJA}[1]{\texorpdfstring{\begin{japanese}#1\end{japanese}}{#1}}
\newcommand{\switchcolumnde}{\switchcolumn}
\newcommand{\switchcolumnla}{}
\newcommand{\switchcolumnja}{\switchcolumn}
\usepackage[de-ja-grc]{attisch-layout}

\begin{document}
\frontmatter
\renewcommand\thepart{\Alph{part}}

% \listoftodos
% \cleardoublepage

\title{{\fontfamily\rmdefault\selectfont アッティカ方言で話しませんか?\\Sprechen Sie Attisch?}\\
  \rule[1.2ex]{0.2\columnwidth}{1pt}\\
  \normalsize
  最良のアッティカの著者に基づく\\
  \kenten{古代ギリシャ語での}\\
  現代（訳注：19世紀末ドイツ）日常会話\\
  \large
  \medskip
  日本語訳\\
  \medskip
  \begin{tabular}{@{第}c@{部\quad}l}
  D&交際にて\\
  E&恋の幸せと恋の苦しみ\\
  F&家の中で\\
  G&政治生活より
  \end{tabular}
}


\author{
  \textlatin{\textbf{E. Joannides博士} 著}\\
  \textlatin{\textbf{ざぎん @na4zagin3} 訳}\\
  {\tiny\begin{minipage}{\columnwidth}
      \begin{flushright}
        \textlatin{--- --- Ridentem discere Graeca\\
          Quid vetat? --- ---}
      \end{flushright}
    \end{minipage}}\smallskip{}
  \\
  \vspace{\fill}
  \rule[1.2ex]{0.2\columnwidth}{1pt}}


\date{
  発行 ヒュアリニオス\\
  コミックマーケット108（2026年夏）\\
  原著 \textbf{Leipzig, 1889.} \emph{C.\ A.\ Koch's Verlag.}
  (J. Sengbusch)}

\pagestyle{empty}
\maketitle
\pagebreak{}
{
\setlength{\parindent}{1\zw}
\setlength{\hangindent}{0pt}
\section*{まえがき}
本書は、%
第A部から第J部、並びに文法要項と単語集からなる原書1889年刊行版%
\url{https://books.google.co.jp/books/about/Sprechen_Sie_Attisch.html?id=swUZAAAAYAAJ&redir_esc=y}%
のうち、本文第D部から第G部までのドイツ語部分を、ギリシャ語部分も参考にして和訳したものです。

Lua\LaTeX でも\XeLaTeX 同様にUniFrakturで髭文字を組むことに成功したので、今回よりドイツ語原文も併記することにしました。

訳の際には既存の英訳%
\url{https://paulatinygriego.files.wordpress.com/2019/02/greekenglishphrasebook.pdf}%
や、アリストパネースの和訳『ギリシア喜劇全集第1巻〜第4巻』（岩波書店）%
も参照しました。

オープンソースプロジェクトとして翻訳をしていきたいと思っているので、
誤りや改善案等ございましたら、是非とも%
\url{https://github.com/na4zagin3/Sprechen-Sie-Attisch}%
にissueを立てるなり、pull requestを送るなりして下さると、幸甚です。

ギリシャ語学習の一助になれば幸いです。一緒に古典語学習ブームを盛り上げていきましょう！

\subsection*{本文について}
本文は3段、左からドイツ語原文、日本語訳、ギリシャ語原文となっています。

\subsection*{文字について}
ドイツ語本文は原文通り亀甲文字で組んでいます。

以下が文字一覧です。
\textgerman{Aa Bb Cc Dd Ee Ff Gg Hh Ii Jj Kk Ll Mm Nn Oo Pp Qq Rr Sss Tt Uu Vv
Ww Xx Yy Zz Ää Öö Üü ß}

以下のように形で分類すると分かりやすいでしょう。
\begin{latin}%
  A \textgerman[spelling=old,babelshorthands=true]{Aa}\quad
  U \textgerman[spelling=old,babelshorthands=true]{Uu;}\qquad
  C \textgerman[spelling=old,babelshorthands=true]{Cc}\quad
  E \textgerman[spelling=old,babelshorthands=true]{Ee}\quad
  S \textgerman[spelling=old,babelshorthands=true]{Sss}\quad
  G \textgerman[spelling=old,babelshorthands=true]{Gg;}\qquad
  K \textgerman[spelling=old,babelshorthands=true]{Kk}\quad
  H \textgerman[spelling=old,babelshorthands=true]{Hh;}\qquad
  N \textgerman[spelling=old,babelshorthands=true]{Nn}\quad
  R \textgerman[spelling=old,babelshorthands=true]{Rr}\quad
  X \textgerman[spelling=old,babelshorthands=true]{Xx;}\qquad
  M \textgerman[spelling=old,babelshorthands=true]{Mm}\quad
  W \textgerman[spelling=old,babelshorthands=true]{Ww}\quad
  V \textgerman[spelling=old,babelshorthands=true]{Vv}\quad
  B \textgerman[spelling=old,babelshorthands=true]{Bb}\quad
  Y \textgerman[spelling=old,babelshorthands=true]{Yy;}\qquad
  O \textgerman[spelling=old,babelshorthands=true]{Oo}\quad
  Q \textgerman[spelling=old,babelshorthands=true]{Qq}\quad
  P \textgerman[spelling=old,babelshorthands=true]{Pp}\quad
  D \textgerman[spelling=old,babelshorthands=true]{Dd;}\qquad
  T \textgerman[spelling=old,babelshorthands=true]{Tt}\quad
  L \textgerman[spelling=old,babelshorthands=true]{Ll;}\qquad
  IJ \textgerman[spelling=old,babelshorthands=true]{Iij}\quad
  F \textgerman[spelling=old,babelshorthands=true]{Ff;}\qquad
  Z \textgerman[spelling=old,babelshorthands=true]{Zz}\quad
  ß \textgerman[spelling=old,babelshorthands=true]{ß}
\end{latin}%

\pagestyle{plain}
\clearpage{}
%\input{legend-ja.tex}
% \begin{quote}
%   
% \end{quote}
\pagebreak{}\gappto\captionsjapanese{%
\renewcommand{\contentsname}{目次}%
\renewcommand{\prepartname}{第}
\renewcommand{\postpartname}{部}
} 
\captionsjapanese
\renewcommand{\contentsname}{目次}
\renewcommand{\prepartname}{第}
\renewcommand{\postpartname}{部}

\begin{quote}
    \setlength{\columnseprule}{0.5pt}\begin{multicols}{2}\tableofcontents{}

\end{multicols}

\pagebreak{}\mainmatter
\end{quote}

% \begin{paracol}{2}
% 1
\section*{Kleine Regeln und Beobachtungen}
\switchcolumn
\section*{些細な規則と注意\addcontentsline{toc}{section}{\emph{些細な規則と注意}}}
\switchcolumn*


% 2
\setlist[enumerate]{parsep=0pt,itemsep=0pt,partopsep=0pt,topsep=0pt}
\begin{enumerate}[leftmargin=0pt,rightmargin=0pt,listparindent =1cm,labelindent=1cm,labelsep=1ex,labelwidth={*},itemindent={*},align=left]
% 3
\begin{german}
\item Nichts erleichtert es so sehr, eine Sprache zu beherrschen, als wenn
man ihre \emph{Schwächen} erspäht. Erst wenn wir ermittelt haben,
was einer Sprache fehlt, verstehen wir recht, warum sie gerade diese
oder jene Wendung vorzieht, diese oder jene Verbindung von Begriffen
liebt, warum sie in dieser oder jener Weise von der Aus\textcompwordmark{}drucks\textcompwordmark{}weise
unserer eigenen Sprache abweicht. Wir erfassen als\textcompwordmark{}dann
ein gutes Theil von ihrem \quotedblbase Geiste``, wie man den Inbegriff
ihrer Besonderheiten so gern nennt.
\end{german}
\switchcolumn

\begin{japanese}
\item その言語の習得を容易たらしめるものとして、言語の\emph{弱点}を見抜くことほどのものはない。ある言語にどこが欠けているかが分かって初めて、ある言い回しを選んだり、ある語の組み合わせを好む理由、そして我々の母語の表現法と異なる理由が真に理解できるのだ。その言語の\emph{精神}——即ち特殊性の総体——の肝が掴めるようになる。
\end{japanese}
\switchcolumn*

% 4
\begin{german}
Eine bemerkens\textcompwordmark{}werthe Schwäche der griechischen
Sprache nun ist es, daß ihr bei allem Formenreichthum doch ein bequem
zu verwendendes \emph{\,Passivum fehlt}. Die Übereinstimmung eines
großen Theiles der passiven Formen mit den medialen erschwert ihre
Anwendung, weil Deutlichkeit das erste Gesets der Sprache ist, und
vielen Zeitwörtern fehlen überdies die allein dem Passivum eigenen
Formen.
\end{german}
\switchcolumn

\begin{japanese}
ギリシア語の顕著な弱点は、形態の豊かさにもかかわらず\emph{使いやすい受動態形が欠けている}ことだ。受動態形の多くが中動態形と同形であることは——明確さは言語の第一原理である故——応用を難しくしている。どの道、多くの動詞には受動専用の形がそもそも存在しないのだ。\todo{この節ここまで校了}
\end{japanese}
\switchcolumn*

\begin{german}
Um die eigenthümliche Färbung der griechischen Sprache nachzuahmen,
hat man daher zu allererst Folgendes zu beachten:
\end{german}
\switchcolumn

\begin{japanese}
ギリシア語特有の色合いをまねるには、まず何よりも次の点に注意しなければならない。
\end{japanese}
\switchcolumn*

\begin{german}
\emph{Man meide thunlichst die den medialen gleichlautenden passiven
Formen und achte darauf, wie der Grieche diese zu ersetzen pflegt.}
\end{german}
\switchcolumn

\begin{japanese}
\emph{中動態と同形の受動形はできるだけ避け、ギリシア人がそれをどのように置き換えるのが常であったかに注意すること。}
\end{japanese}
\switchcolumn*

\begin{german}
Nur die durch den Zusammenhang sofort als solche erkennbaren und
gewisse in häufigen Gebrauch gekommene Passiva der bezeichneten Art
sind unbedenklich anzuwenden.
\end{german}
\switchcolumn

\begin{japanese}
この種の受動形のうち安心して用いられるのは、文脈からただちに受動と分かるもの、および頻用されて慣用化した一定のものに限られる。
\end{japanese}
\switchcolumn*

\begin{german}
Umschreibungen des Passivums geschehen.
\end{german}
\switchcolumn

\begin{japanese}
受動の迂言は次のように行われる。
\end{japanese}
\switchcolumn*

\begin{enumerate}
\begin{german}
\item durch active Verba, z.\,B.
\end{german}
\switchcolumn

\begin{japanese}
能動動詞による。例えば、
\end{japanese}
\switchcolumn*

\begin{german}
\begin{quote}
belehrt werden \textgreek[variant=ancient]{μανθάνειν,}\\
gerühmt werden \textgreek[variant=ancient]{εὐδοκιμεῖν,}\\
geplagt werden \textgreek[variant=ancient]{κάμνειν,}\\
vor Gericht gestellt werden \textgreek[variant=ancient]{εἰσιέναι εἰς
δικαστήριον,}\\
verklagt werden \textgreek[variant=ancient]{φεύγειν,}\\
gehalten werden für \ldots{} \textgreek[variant=ancient]{δοκεῖν,}\\
es wird mir etwas zugefügt \textgreek[variant=ancient]{πάσχω τι,}\\
vertrieben werden \textgreek[variant=ancient]{ἐκπίπτειν,}\\
einer Sache beraubt werden \textgreek[variant=ancient]{ἀπολλύναι τι},\\
getödtet werden \textgreek[variant=ancient]{ἀποθνήσκειν},\\
sie wurden vertrieben \textgreek[variant=ancient]{ἀνέστησαν},\\
es wurde mir geantwortet \textgreek[variant=ancient]{ἤκουσα},\\
es wird mir Gutes erwiesen \textgreek[variant=ancient]{εὖ πάσχω},\\
ich ward durch's Loos gewählt \textgreek[variant=ancient]{ἔλαχον},\\
ich ward freigesprochen \textgreek[variant=ancient]{ἀπέφυγον},\\
ich ward geschmäht \textgreek[variant=ancient]{κακῶς ἤκουσα},\\
ich ward (von Mitleid) ergriffen \textgreek[variant=ancient]{(ἔλεός)
με εἰσῄει}.
\end{quote}
\end{german}
\switchcolumn

\begin{japanese}
\begin{quote}
教えられる \textgreek[variant=ancient]{μανθάνειν,}\\
称賛される \textgreek[variant=ancient]{εὐδοκιμεῖν,}\\
悩まされる \textgreek[variant=ancient]{κάμνειν,}\\
法廷に立たされる \textgreek[variant=ancient]{εἰσιέναι εἰς
δικαστήριον,}\\
訴えられる \textgreek[variant=ancient]{φεύγειν,}\\
……と見なされる \textgreek[variant=ancient]{δοκεῖν,}\\
私に何かが加えられる \textgreek[variant=ancient]{πάσχω τι,}\\
追放される \textgreek[variant=ancient]{ἐκπίπτειν,}\\
ある物を奪われる \textgreek[variant=ancient]{ἀπολλύναι τι},\\
殺される \textgreek[variant=ancient]{ἀποθνήσκειν},\\
彼らは追い立てられた \textgreek[variant=ancient]{ἀνέστησαν},\\
私に返答があった \textgreek[variant=ancient]{ἤκουσα},\\
私はよい扱いを受ける \textgreek[variant=ancient]{εὖ πάσχω},\\
私はくじで選ばれた \textgreek[variant=ancient]{ἔλαχον},\\
私は無罪放免となった \textgreek[variant=ancient]{ἀπέφυγον},\\
私は悪しざまに言われた \textgreek[variant=ancient]{κακῶς ἤκουσα},\\
私は（憐れみに）襲われた \textgreek[variant=ancient]{(ἔλεός)
με εἰσῄει}.
\end{quote}
\end{japanese}
\switchcolumn*

\begin{german}
\item vielfach durch \textgreek[variant=ancient]{γίγνεσθαι;} es steht für
gemacht, veranstaltet, bewerkstelligt werden, übertragen, verliehen,
erkauft, erworben werden, verübt w., gefeiert w. (von Festen), geboren
w. und andere Passiva.
\end{german}
\switchcolumn

\begin{japanese}
多くの場合、\textgreek[variant=ancient]{γίγνεσθαι}による。
これは「作られる」「催される」「成し遂げられる」「移される」「授与される」
「買い取られる」「獲得される」「（悪事が）犯される」「（祭りが）祝われる」
「生まれる」などの受動に代わって用いられる。
\end{japanese}
\switchcolumn*

\begin{german}
\item durch Substantiva mit Verben, z.\,B.
\end{german}
\switchcolumn

\begin{japanese}
名詞と動詞による。例えば、
\end{japanese}
\switchcolumn*

\begin{german}
\begin{quote}
gelobt werden \textgreek[variant=ancient]{ἔπαινον ἔχειν,}\\
es wird (viel) gesprochen \textgreek[variant=ancient]{λόγος ἐστὶ (πολύς),}\\
bestraft werden \textgreek[variant=ancient]{δίκην διδόναι,}\\
es wird gezürnt u. \textgreek[variant=ancient]{ὀργὴ γίγνεται} dgl.\ mehr;
\end{quote}
\end{german}
\switchcolumn

\begin{japanese}
 \begin{quote}
褒められる \textgreek[variant=ancient]{ἔπαινον ἔχειν,}\\
（多く）語られる \textgreek[variant=ancient]{λόγος ἐστὶ (πολύς),}\\
罰せられる \textgreek[variant=ancient]{δίκην διδόναι,}\\
怒られる、など \textgreek[variant=ancient]{ὀργὴ γίγνεται} ほか。
\end{quote}
\end{japanese}
\switchcolumn*

\begin{german}
\item durch Adjektiva mit \textgreek[variant=ancient]{εἶναι,} z.\,B.
\end{german}
\switchcolumn

\begin{japanese}
形容詞＋\textgreek[variant=ancient]{εἶναι}による。例えば、
\end{japanese}
\switchcolumn*

\begin{german}
\begin{quote}
gesehen werden \textgreek[variant=ancient]{καταφανῆ εἰναι},\\
es wird dir nicht geglaubt \textgreek[variant=ancient]{ἄπιστος εἶ}
u. dgl. mehr.
\end{quote}
\end{german}
\switchcolumn

\begin{japanese}
 \begin{quote}
見られる \textgreek[variant=ancient]{καταφανῆ εἰναι},\\
あなたは信じてもらえない \textgreek[variant=ancient]{ἄπιστος εἶ}
など。
\end{quote}
\end{japanese}
\switchcolumn*


\end{enumerate}
\begin{german}
\item Im Griechischen fehlt die Genauigkeit in der Bezeichnung des Objectes,
wie sie den modernen Sprachen eigen ist. Die letzteren setzen, wenn
zwei verbundene Verba das gleiche Object in verschiedenem Casus erfordern,
zum zweiten Verbum anstatt der Wiederholung des Nomens das persönliche
Pronomen (seiner, ihm, ihn, ihrer, ihr, sie, es, ihnen) als Object,
\emph{der Grieche läßt die Stelle des gemeinsamen Objectes beim zweiten
Verbum unbezeichnet, gleichviel in welchem Casus es stehen müßte.}
\end{german}
\switchcolumn

\begin{japanese}
ギリシア語には、近代諸語に見られるような目的語の指示の精密さが欠けている。近代語では、
二つの動詞が結びついて同一の目的語を異なる格で要求する場合、二つ目の動詞には名詞を繰り返さず、
人称代名詞（彼の、彼に、彼を、彼女の、彼女に、彼女を、それ、それら、彼らに等）を目的語として置く。
\emph{しかしギリシア語では、二つ目の動詞の位置に共通の目的語を示さない。必要な格が何であろうと同じである。}
\end{japanese}
\switchcolumn*

\begin{german}
Das dem französischen \textfrench{en} entsprechende Object (welchen,
welche, welches) wird im Griechischen nicht aus\textcompwordmark{}gedrückt,
z.\,B.: Sie werden das Gold aus Lydien holen lassen müssen, wenn
sie welches haben wollen \textgreek[variant=ancient]{ἐκ Λυδίας μεταστέλλεσθαι
τὸ χρυσίον δεήσει αὐτοὺς, ἢν ἐπιθυμήσωσιν.}
\end{german}
\switchcolumn

\begin{japanese}
フランス語の \textfrench{en} に相当する目的語（「それを／どれを」）は、
ギリシア語では表されない。例えば「彼らがそれを欲するなら、リュディアから金を取り寄せねばならない」
\textgreek[variant=ancient]{ἐκ Λυδίας μεταστέλλεσθαι τὸ χρυσίον δεήσει αὐτοὺς, ἢν ἐπιθυμήσωσιν.}
\end{japanese}
\switchcolumn*

\begin{german}
\item Dem Griechen fehlt, wie dem Lateiner, das Mittel zur Hervorhebung
einzelner Satztheile, welches unsere Sprache, ähnlich anderen modernen
Sprachen, darin besitzt, daß sie den hervorzuhebenden Begriff zum
Prädivcte eines neuen Satzes meist mit dem unpersönlichen Subject
es macht, während die übrigen Satztheile in einem abhängigen Satze
vermittelst eines Relativs oder einer Conjunction angefügt werden.
\emph{Im Griechischen muß die der Hervorhebung eines Begriffes dienende
Zerlegung eines Satzes in zwei unterbleiben,} z.\,B.:  Es ist derselbe,
der dies sagt \textgreek[variant=ancient]{ὁ αὐτὸς ταῦτα λέγει}. Wer
ist der Mann, den du rufst? \textgreek[variant=ancient]{τίνα τὸν ἄνδρα
καλεῖς;} Ist es wahr, daß du das gethan hast? \textgreek[variant=ancient]{ἆρ᾽
ἀληθῶς τοῦτ᾽ ἐποίησας;} Wie ist es möglich, daß\ldots{} \textgreek[variant=ancient]{πῶς\ldots{};}
wie kommt es, daß\ldots{} \textgreek[variant=ancient]{πῶς\ldots{}; }
\end{german}
\switchcolumn

\begin{japanese}
ギリシア語にもラテン語にも、近代諸語に見られる次のような強調構文はない。
ドイツ語では、他の近代諸語と同じく、強調したい概念を、多くは非人称主語「それ」を伴う新たな文の述語とし、
残りの文成分を関係詞や接続詞で導かれる従属節として添える。これに対し、
ギリシア語では\emph{強調のために文を二つに分けることができない}。例えば、
「それを言うのは彼だ」\textgreek[variant=ancient]{ὁ αὐτὸς ταῦτα λέγει}。
「君が呼んでいる男は誰だ」\textgreek[variant=ancient]{τίνα τὸν ἄνδρα καλεῖς;}。
「君がそれをしたというのは本当か」\textgreek[variant=ancient]{ἆρ᾽ ἀληθῶς τοῦτ᾽ ἐποίησας;}。
「どうしてそんなことが可能なのか」\textgreek[variant=ancient]{πῶς\ldots{};}、
「どうしてそうなったのか」\textgreek[variant=ancient]{πῶς\ldots{};}。
\end{japanese}
\switchcolumn*

\begin{german}
\item Coordinirte Sätze und coordinirte Satztheile kann der Grieche nicht
unverbunden lassen. Asyndetisches Nebeneinanderstellen von Satztheilen
kommt nur selten und zwar als Aus\textcompwordmark{}druck lebhafter
Erregung zur Anwendung.
\end{german}
\switchcolumn

\begin{japanese}
ギリシア語では、並列する文や並列する句を結びつけないままにしておくことができない。
接続詞省略によって並べる用法はまれで、激しい感情の表現として用いられる場合に限られる。
\end{japanese}
\switchcolumn*

\begin{german}
In ununterbrochener Rede ist \emph{jeder neue Satz} durch eine passende
Conjunction (\textgreek[variant=ancient]{δέ, καί οὖν, γάρ} etc.) an
das Voraus\textcompwordmark{}gehende \emph{anzuschließen.}
\end{german}
\switchcolumn

\begin{japanese}
途切れない叙述の中では、\emph{新しい文は必ず}適切な接続詞
（\textgreek[variant=ancient]{δέ, καί οὖν, γάρ} など）で前文に\emph{結び付ける}必要がある。
\end{japanese}
\switchcolumn*

\begin{german}
Der Lernende ist davor zu warnen, \textgreek[variant=ancient]{μέν}
für eine diese Verbin- dung mit dem \emph{Voraus\textcompwordmark{}gehenden}
ersetzende Conjunction zu halten, da es nur zum Hinweis auf das \emph{Folgende}
dient.
\end{german}
\switchcolumn

\begin{japanese}
学習者は、\textgreek[variant=ancient]{μέν} を前文との結び付きに代わる接続詞だと考えてはならない。
それは\emph{後続}への指示にしか役立たないからである。
\end{japanese}
\switchcolumn*

\begin{german}
Anfügung ohne Bindewort ist in ununterbrochener Rede nur gestattet:
\end{german}
\switchcolumn

\begin{japanese}
接続語なしの挿入は、途切れない叙述では次の場合にのみ許される。
\end{japanese}
\switchcolumn*

\begin{enumerate}
\begin{german}
\item an den Stellen, wo wir im Deutschen den Doppelpunkt als Interpunctions\textcompwordmark{}zeichen
setzen;
\end{german}
\switchcolumn

\begin{japanese}
ドイツ語でコロンを置く場所。
\end{japanese}
\switchcolumn*

\begin{german}
\item wenn der neue Satz mit stark betontem Demonstrativum oder
\end{german}
\switchcolumn

\begin{japanese}
新しい文が強く強調された指示詞で始まるとき。
\end{japanese}
\switchcolumn*

\begin{german}
\item wenn der neue Satz mit \textgreek[variant=ancient]{εἶτα} (= und dann)
oder \textgreek[variant=ancient]{ἔπειτα} beginnt;
\end{german}
\switchcolumn

\begin{japanese}
新しい文が \textgreek[variant=ancient]{εἶτα}（＝「そして」）や
\textgreek[variant=ancient]{ἔπειτα} で始まるとき。
\end{japanese}
\switchcolumn*

\begin{german}
\item wo wir im Deutschen mit \emph{\quotedblbase nicht aber``} fortsahren\footnote{\textjapanese{原文はfortsahren（誤記。fortfahrenのこと）。}};
es steht dann häufig bloßes \textgreek[variant=ancient]{οὐ} (beziehentlich
\textgreek[variant=ancient]{μή}\,), (weil \textgreek[variant=ancient]{οὐ}
mit \textgreek[variant=ancient]{δέ} \quotedblbase und nicht`` oder
\quotedblbase nicht einmal`` bedeutet), oft jedoch auch \textgreek[variant=ancient]{οὐ
μέντοι.}
\end{german}
\switchcolumn

\begin{japanese}
ドイツ語で「しかし〜ではない」と続ける場所。そこでしばしば単独の
\textgreek[variant=ancient]{οὐ}（あるいは \textgreek[variant=ancient]{μή}）が置かれる
（\textgreek[variant=ancient]{οὐ} が \textgreek[variant=ancient]{δέ} と結び付くと
「しかし〜ではない」「〜すらない」の意になるため）。また
\textgreek[variant=ancient]{οὐ μέντοι} が用いられることも多い。
\end{japanese}
\switchcolumn*

\end{enumerate}
\begin{german}
\item Man merke: Nun so ... denn = \textgreek[variant=ancient]{ἀλλά,}
\end{german}
\switchcolumn

\begin{japanese}
注意：さて、では＝\textgreek[variant=ancient]{ἀλλά,}
\end{japanese}
\switchcolumn*

\bgroup\addtolength{\leftskip}{\parindent}\addtolength{\leftskip}{3em}\setlength{\parindent}{-1em}\setlength{\arraycolsep}{0pt}%
\begin{german}
o dann ... = \textgreek[variant=ancient]{ἄρα,}
\end{german}
\switchcolumn

\begin{japanese}
おお、それなら……＝\textgreek[variant=ancient]{ἄρα,}
\end{japanese}
\switchcolumn*

\begin{german}
da kam, da sagte = \textgreek[variant=ancient]{καὶ ἦλθε, καὶ εἰπεν,}
\end{german}
\switchcolumn

\begin{japanese}
そこで来て、そこで言った＝\textgreek[variant=ancient]{καὶ ἦλθε, καὶ εἰπεν,}
\end{japanese}
\switchcolumn*

\begin{german}
jedoch = \textgreek[variant=ancient]{μέντοι,}
\end{german}
\switchcolumn

\begin{japanese}
しかしながら＝\textgreek[variant=ancient]{μέντοι,}
\end{japanese}
\switchcolumn*

\begin{german}
denn sonst \ldots{} =\textgreek[variant=ancient]{γάρ,}
\end{german}
\switchcolumn

\begin{japanese}
さもなくば……＝\textgreek[variant=ancient]{γάρ,}
\end{japanese}
\switchcolumn*

\begin{german}
denn (folgernd), z.B. höre \emph{denn, so} ward er \emph{denn} ..
= \textgreek[variant=ancient]{δή,}
\end{german}
\switchcolumn

\begin{japanese}
（推論の）それで、だから。例：\emph{それで}彼は……＝\textgreek[variant=ancient]{δή,}
\end{japanese}
\switchcolumn*

\begin{german}
doch wohl (ohne Zweifel) = \textgreek[variant=ancient]{δήπου,}
\end{german}
\switchcolumn

\begin{japanese}
きっと（疑いなく）＝\textgreek[variant=ancient]{δήπου,}
\end{japanese}
\switchcolumn*

\begin{german}
und schon = \textgreek[variant=ancient]{καὶ δή} (\textgreek[variant=ancient]{δή}
= \textgreek[variant=ancient]{ἤδη}), vgl. \textgreek[variant=ancient]{πάλαι
δή} schon längst, \textgreek[variant=ancient]{νῦν δή} jetzt eben,
\end{german}
\switchcolumn

\begin{japanese}
そしてすでに＝\textgreek[variant=ancient]{καὶ δή}
（\textgreek[variant=ancient]{δή}＝\textgreek[variant=ancient]{ἤδη}）。
例：\textgreek[variant=ancient]{πάλαι δή}「すでにはるか以前」、
\textgreek[variant=ancient]{νῦν δή}「まさに今」。
\end{japanese}
\switchcolumn*

\begin{german}
wohl aber = \textgreek[variant=ancient]{δὲ,}
\end{german}
\switchcolumn

\begin{japanese}
しかし一方で＝\textgreek[variant=ancient]{δὲ,}
\end{japanese}
\switchcolumn*

\begin{german}
\begin{tabular}{lc}
dann erst & \ldelim\}{2}{1em}[]\tabularnewline
erst dann & \tabularnewline
\end{tabular} =\footnote{\textjapanese{原文は等号を欠く}} \textgreek[variant=ancient]{οὕτω
δή,}
\end{german}
\switchcolumn

\begin{japanese}
それで初めて／その時になって＝\textgreek[variant=ancient]{οὕτω δή,}
\end{japanese}
\switchcolumn*

\begin{german}
\ldots{} allerdings = \textgreek[variant=ancient]{\ldots{} μήν,}
\end{german}
\switchcolumn

\begin{japanese}
……もっとも＝\textgreek[variant=ancient]{\ldots{} μήν,}
\end{japanese}
\switchcolumn*

\begin{german}
indessen \ldots{} = \textgreek[variant=ancient]{οὐ μὴν ἀλλά,}
\end{german}
\switchcolumn

\begin{japanese}
しかし一方で……＝\textgreek[variant=ancient]{οὐ μὴν ἀλλά,}
\end{japanese}
\switchcolumn*

\begin{german}
wahrscheinlich (adv.) =\footnote{\textjapanese{原文は等号を欠く}}%
\textgreek[variant=ancient]{ἦ που \ldots{}}
\end{german}
\switchcolumn

\begin{japanese}
おそらく（副詞）＝\textgreek[variant=ancient]{ἦ που \ldots{}}
\end{japanese}
\switchcolumn*

\begin{german}
oder (nach Negationen) = \textgreek[variant=ancient]{οὐδέ, μνδέ,}
\end{german}
\switchcolumn

\begin{japanese}
（否定後の）「あるいは」＝\textgreek[variant=ancient]{οὐδέ, μνδέ,}
\end{japanese}
\switchcolumn*

\begin{german}
doch (lat. \textlatin{quaeso}) = \textgreek[variant=ancient]{δῆτα,}
\end{german}
\switchcolumn

\begin{japanese}
どうか／お願い（ラテン語 \textlatin{quaeso}）＝\textgreek[variant=ancient]{δῆτα,}
\end{japanese}
\switchcolumn*

\begin{german}
nicht sowohl \textgreek[variant=ancient]{\ldots{}} als vielmehr =
\footnote{\textjapanese{原文は波括弧が左右反転している}}%
\end{german}
\switchcolumn

\begin{japanese}
\textgreek[variant=ancient]{\ldots{}} というよりむしろ……＝
\end{japanese}
\switchcolumn*

\begin{german}
\begin{tabular}{ll}
\ldelim\{{2}{1em}[] & \begin{greek}[variant=ancient]%
οὐ τοσοῦτον ὅσον \ldots{}\end{greek}%
\tabularnewline
 & \begin{greek}[variant=ancient]%
οὐ τὸ πλέον \ldots{} ἀλλὰ \ldots{}\end{greek}%
\tabularnewline
\end{tabular}
\end{german}
\switchcolumn

\begin{japanese}
\begin{tabular}{ll}
\ldelim\{{2}{1em}[] & \begin{greek}[variant=ancient]%
……ほどではなく……\end{greek}%
\tabularnewline
 & \begin{greek}[variant=ancient]%
……というよりむしろ……\end{greek}%
\tabularnewline
\end{tabular}
\end{japanese}
\switchcolumn*
\egroup

\begin{german}
Aus der Thatsache, daß \quotedblbase o dann ...`` sich
überall passend durch \textgreek[variant=ancient]{ἄρα} geben läßt,
folgt noch keines\textcompwordmark{}wegs, daß umgekehrt \textgreek[variant=ancient]{ἄρα}
sich überall passend durch \quotedblbase o dann \ldots{}`` übersetzen
lasse.
\end{german}
\switchcolumn

\begin{japanese}
「おお、それなら……」がどこでも \textgreek[variant=ancient]{ἄρα} で適切に言えるからといって、
逆に \textgreek[variant=ancient]{ἄρα} をどこでも「おお、それなら……」と訳せるわけではない。
\end{japanese}
\switchcolumn*

\begin{german}
\item \emph{Großes} Glück \textgreek[variant=ancient]{πολλὴ εὐδαιμονία.}
\end{german}
\switchcolumn

\begin{japanese}
\emph{大きな}幸運 \textgreek[variant=ancient]{πολλὴ εὐδαιμονία.}
\end{japanese}
\switchcolumn*

\begin{continuousitemline}


\begin{german}
Großes Mißgeschick \textgreek[variant=ancient]{πολλὴ δυστυχία.}
\end{german}
\switchcolumn

\begin{japanese}
大きな不運 \textgreek[variant=ancient]{πολλὴ δυστυχία.}
\end{japanese}
\switchcolumn*

\begin{german}
Großer Überfluß \textgreek[variant=ancient]{πολλὴ ἀφθονία.}
\end{german}
\switchcolumn

\begin{japanese}
大いなる豊富 \textgreek[variant=ancient]{πολλὴ ἀφθονία.}
\end{japanese}
\switchcolumn*

\begin{german}
Große Thorheit \textgreek[variant=ancient]{πολλὴ μωρία.}
\end{german}
\switchcolumn

\begin{japanese}
大きな愚かさ \textgreek[variant=ancient]{πολλὴ μωρία.}
\end{japanese}
\switchcolumn*

\begin{german}
Große Unwissenheit \textgreek[variant=ancient]{πολλὴ ἀμαθία.}
\end{german}
\switchcolumn

\begin{japanese}
大きな無知 \textgreek[variant=ancient]{πολλὴ ἀμαθία.}
\end{japanese}
\switchcolumn*

\begin{german}
Große Unvernunft \textgreek[variant=ancient]{πολλὴ ἀλογία.}
\end{german}
\switchcolumn

\begin{japanese}
大きな不合理 \textgreek[variant=ancient]{πολλὴ ἀλογία.}
\end{japanese}
\switchcolumn*

\begin{german}
Große Geschäftigkeit \textgreek[variant=ancient]{πολλὴ πραγματεία.}
\end{german}
\switchcolumn

\begin{japanese}
大いなる多忙 \textgreek[variant=ancient]{πολλὴ πραγματεία.}
\end{japanese}
\switchcolumn*

\begin{german}
Sehr große Muthlosigkeit \textgreek[variant=ancient]{πλείστη ἀθυμία.}
\end{german}
\switchcolumn

\begin{japanese}
非常に大きな気後れ \textgreek[variant=ancient]{πλείστη ἀθυμία.}
\end{japanese}
\switchcolumn*
\end{continuousitemline}

\begin{german}
\item %
\begin{tabular}[t]{lc}
So ein trefflicher & \rdelim\}{5}{1em}[{\textgreek[variant=ancient]{τοιοῦτος.}}]\tabularnewline
So ein abscheulicher & \tabularnewline
So ein erfahrener & \tabularnewline
So ein beschränkter & \tabularnewline
So ein gefährlicher & \tabularnewline
\end{tabular}
\end{german}
\switchcolumn

\begin{japanese}
こんな立派な／こんなひどい／こんな経験豊かな／こんな狭量な／こんな危険な
\end{japanese}
\switchcolumn*

\begin{continuousitemline}%
\begin{german}
\qquad{}u. s. w.
\end{german}
\switchcolumn

\begin{japanese}
など
\end{japanese}
\switchcolumn*

\begin{german}
\begin{tabular}[t]{lc}
So ein trefflicher & \rdelim\}{5}{1em}[{ \textgreek[variant=ancient]{τοιοῦτος.}}]\tabularnewline
So ein abscheulicher & \tabularnewline
So ein erfahrener & \tabularnewline
So ein beschränkter & \tabularnewline
So ein gefährlicher & \tabularnewline
\end{tabular}
\end{german}
\switchcolumn

\begin{japanese}
こんな立派な／こんなひどい／こんな経験豊かな／こんな狭量な／こんな危険な
\end{japanese}
\switchcolumn*

\begin{german}
\qquad{}u. s. w.
\end{german}
\switchcolumn

\begin{japanese}
など
\end{japanese}
\switchcolumn*

\begin{german}
\begin{tabular}[t]{lc}
So Verwerfliches & \rdelim\}{2}{1em}[{ \textgreek[variant=ancient]{τοιαῦτα.}}]\tabularnewline
So Löbliches & \tabularnewline
\end{tabular}
\end{german}
\switchcolumn

\begin{japanese}
これほど非難すべきこと／これほど賞賛すべきこと
\end{japanese}
\switchcolumn*

\begin{german}
\qquad{}u. s. w.
\end{german}
\switchcolumn

\begin{japanese}
など
\end{japanese}
\switchcolumn*

\begin{german}
\begin{tabular}[t]{lc}
es klingt schön & \rdelim\}{3}{1em}[{ \textgreek[variant=ancient]{ἡδύ ἐστιν.}}]\tabularnewline
es schmeckt gut & \tabularnewline
es riecht gut & \tabularnewline
\end{tabular}
\end{german}
\switchcolumn

\begin{japanese}
響きが美しい／味がよい／匂いがよい
\end{japanese}
\switchcolumn*

\begin{german}
\begin{tabular}[t]{lc}
(jetzt) so spät & \rdelim\}{2}{1em}[{ \textgreek[variant=ancient]{τηνικάδε.}}]\tabularnewline
(jetzt) so früh & \tabularnewline
\end{tabular}
\end{german}
\switchcolumn

\begin{japanese}
（いま）こんなに遅く／（いま）こんなに早く
\end{japanese}
\switchcolumn*

\begin{german}
\quad{}Der gewöhnliche Aus\textcompwordmark{}druck für
\end{german}
\switchcolumn

\begin{japanese}
……の普通の表現は
\end{japanese}
\switchcolumn*

\begin{german}
\begin{tabular}[t]{lc}
hoffen & \rdelim\}{2}{1em}[\textgerman{ ist \textgreek[variant=ancient]{οἴεσθαι},}]\tabularnewline
fürchten & \tabularnewline
\end{tabular}
\end{german}
\switchcolumn

\begin{japanese}
「望む」「恐れる」の普通の言い方は
\end{japanese}
\switchcolumn*

\begin{german}
\begin{tabular}[t]{lc}
versprechen & \rdelim\}{4}{1em}[\textgerman{ ist \textgreek[variant=ancient]{φάναι.}}]\tabularnewline
brohen\footnote{\textjapanese{原文はbrohen（誤記。drohenのこと）。}} & \tabularnewline
antworten & \tabularnewline
erwidern & \tabularnewline
\end{tabular}
\end{german}
\switchcolumn

\begin{japanese}
「約束する」「脅す」「答える」「言い返す」の普通の言い方は
\end{japanese}
\switchcolumn*

\begin{german}
\ldots{} fuhr er fort, = \textgreek[variant=ancient]{ἔφη.}
\end{german}
\switchcolumn

\begin{japanese}
……と彼は続けた＝\textgreek[variant=ancient]{ἔφη.}
\end{japanese}
\switchcolumn*
\end{continuousitemline}

\begin{german}
\item Ein Freund \textgreek[variant=ancient]{φίλος \emph{τις.}}
\end{german}
\switchcolumn

\begin{japanese}
友人 \textgreek[variant=ancient]{φίλος \emph{τις.}}
\end{japanese}
\switchcolumn*

\begin{german}
\begin{continuousitemline}%
Ein redlicher Freund \textgreek[variant=ancient]{χρηστός
τις \emph{ἄνθρωπος} φίλος.}\par\end{continuousitemline}
\end{german}
\switchcolumn

\begin{japanese}
誠実な友人 \textgreek[variant=ancient]{χρηστός τις \emph{ἄνθρωπος} φίλος.}
\end{japanese}
\switchcolumn*

\begin{german}
\item Unsere 500 Schüler \textgreek[variant=ancient]{\emph{οἱ} ἡμέτεροι
πεντακόσιοι μαθηταί.}
\end{german}
\switchcolumn

\begin{japanese}
私たちの500人の生徒 \textgreek[variant=ancient]{\emph{οἱ} ἡμέτεροι πεντακόσιοι μαθηταί.}
\end{japanese}
\switchcolumn*

\begin{german}
\begin{continuousitemline}Meine drei besten Schüler \textgreek[variant=ancient]{\emph{οἱ}
τρεῖς ἄριστοι τῶν μαθητῶν μου.}\par\end{continuousitemline}
\end{german}
\switchcolumn

\begin{japanese}
私の最良の生徒三人 \textgreek[variant=ancient]{\emph{οἱ} τρεῖς ἄριστοι τῶν μαθητῶν μου.}
\end{japanese}
\switchcolumn*

\begin{german}
\item Ich verlange kein Geld, sondern Zuneigung (Liebe) \textgreek[variant=ancient]{αἰτῶ
\emph{οὐκ} ἀργύριον, ἀλλ᾽ εὔνοιαν.}
\end{german}
\switchcolumn

\begin{japanese}
私は金ではなく好意（愛）を求める \textgreek[variant=ancient]{αἰτῶ \emph{οὐκ} ἀργύριον, ἀλλ᾽ εὔνοιαν.}
\end{japanese}
\switchcolumn*

\begin{german}
\item Ich \emph{habe} gehabt \textgreek[variant=ancient]{εἶχον,} z.\,B.
ich habe ebenfalls diese Klasse einmal gehabt \textgreek[variant=ancient]{κἀγὼ
εἶχον τὴν τάξιν ταύτην ποτέ.} Er \emph{ist} gestern bei mir gewesen
\textgreek[variant=ancient]{παρ᾽ ἐμοὶ χθὲς ἦν.}
\end{german}
\switchcolumn

\begin{japanese}
「持っていた」\textgreek[variant=ancient]{εἶχον}。例えば「私もこのクラスを一度受け持った」
\textgreek[variant=ancient]{κἀγὼ εἶχον τὴν τάξιν ταύτην ποτέ.}
「彼は昨日私の所にいた」\textgreek[variant=ancient]{παρ᾽ ἐμοὶ χθὲς ἦν.}
\end{japanese}
\switchcolumn*

\begin{german}
Das \emph{Perfectum} von \emph{sein} und \emph{haben} und allen ein
Dauer ausdrückenden Verben wird im Griechischen durch das Imperfectum,
bei den übrigen Verben meist durch den Aorist, seltener durch das
Perfectum wiedergegeben. Läßt sich zu dem Verbum ein Adverb der Vergangenheit
(z.\,B. damals) hinzudenken, so steht Aorist; läßt sich ein Adverb
der Gegenwart (z.\,B. nunmehr, bereits) hinzudenken, nur dann steht
Perfectum.
\end{german}
\switchcolumn

\begin{japanese}
「ある」「持つ」および継続を表す動詞の\emph{完了}は、ギリシア語では未完了過去で表される。
それ以外の動詞では多くがアオリスト、まれに完了形で表す。
その動詞に過去の副詞（例：当時）を添えられるならアオリスト、
現在の副詞（例：今や、すでに）を添えられるなら完了形となる。
\end{japanese}
\switchcolumn*

\begin{continuousexamples}%
\begin{german}%
\emph{Hast} du das Geld gefunden? (\textlatin{sc.}
nunmehr) \textgreek[variant=ancient]{ἆρ᾽ εὕρηκας τἀργύριον;}
\end{german}
\switchcolumn

\begin{japanese}
その金は\emph{見つかったか}（今や） \textgreek[variant=ancient]{ἆρ᾽ εὕρηκας τἀργύριον;}
\end{japanese}
\switchcolumn*

\begin{german}
Ja, ich \emph{habe} es gefunden (\textlatin{sc.} nunmehr) \textgreek[variant=ancient]{εὕρηκανὴ
Δία.}
\end{german}
\switchcolumn

\begin{japanese}
はい、\emph{見つけました}（今や） \textgreek[variant=ancient]{εὕρηκανὴ Δία.}
\end{japanese}
\switchcolumn*

\begin{german}
Wo \emph{hast} du es gefunden? (\textlatin{sc.} damals als du es fandest)
\textgreek[variant=ancient]{ποῦ εὗρες;}
\end{german}
\switchcolumn

\begin{japanese}
どこで\emph{見つけたのか}（その時） \textgreek[variant=ancient]{ποῦ εὗρες;}
\end{japanese}
\switchcolumn*

\begin{german}
Ich \emph{habe} es (\textlatin{sc.} damals) in dem Garden\footnote{\textjapanese{原文はGarden（誤記。Gartenのこと）。}} gefunden
\textgreek[variant=ancient]{ἐν τῷ κήπῳ εὗρον. }
\end{german}
\switchcolumn

\begin{japanese}
私は（その時）庭で\emph{見つけた} \textgreek[variant=ancient]{ἐν τῷ κήπῳ εὗρον.}
\end{japanese}
\switchcolumn*
\end{continuousexamples}

\begin{german}
\item Der Infinitiv Aoristi bezeichnet nach den Verben des Sagens und Meinens
die Vergangenheit, z.\,B.
\end{german}
\switchcolumn

\begin{japanese}
アオリスト不定詞は、言う・思う類の動詞に続くとき過去を表す。例えば、
\end{japanese}
\switchcolumn*

\begin{continuousexamples}
\begin{german}
\textgreek[variant=ancient]{φησὶν εὑρεῖν}
er behauptet er \emph{habe} gefunden.
\end{german}
\switchcolumn

\begin{japanese}
\textgreek[variant=ancient]{φησὶν εὑρεῖν}＝彼は\emph{見つけた}と言っている。
\end{japanese}
\switchcolumn*
\end{continuousexamples}

\begin{german}
\item Bedeutet \emph{daß} soviel wie \emph{mache(t)} daß, so wird es durch
\textgreek[variant=ancient]{ὅπως}\footnote{\begin{latin}%
orig. \textgreek[variant=ancient]{οπως}\end{latin}%
} mit dem \textlatin{Indic. Fut.} ausgedrückt.
\end{german}
\switchcolumn

\begin{japanese}
\emph{daß} が「〜させよ／〜になるように」という意味なら、
\textgreek[variant=ancient]{ὅπως}＋未来直説法で表す。
\end{japanese}
\switchcolumn*

\begin{german}
\begin{continuousexamples}Daß es nur kein Mensch erfährt! \textgreek[variant=ancient]{ὅπως
ταῦτα μηδεὶς ἀνθρώπων πεύσεται!}\par\end{continuousexamples}
\end{german}
\switchcolumn

\begin{japanese}
どうか誰にも知られませんように！ \textgreek[variant=ancient]{ὅπως ταῦτα μηδεὶς ἀνθρώπων πεύσεται!}
\end{japanese}
\switchcolumn*

\begin{german}
\item Mit \textgreek[variant=ancient]{ἐξ οὗ} oder \textgreek[variant=ancient]{ἐπεί}
= \emph{seit} verträgt sich kein \textgreek[variant=ancient]{οὐ} oder
\textgreek[variant=ancient]{μή}\footnote{\begin{latin}%
orig. \textgreek[variant=ancient]{μη}\end{latin}%
}\textgreek[variant=ancient]{.}
\end{german}
\switchcolumn

\begin{japanese}
\textgreek[variant=ancient]{ἐξ οὗ} または \textgreek[variant=ancient]{ἐπεί}
（＝「〜以来」）には \textgreek[variant=ancient]{οὐ} や \textgreek[variant=ancient]{μή} を伴わない。
\end{japanese}
\switchcolumn*

\begin{german}
\begin{continuousexamples}Seit wir uns \emph{nicht} gesehen, hat
es viel geregnet: \textgreek[variant=ancient]{ἐξ οὗ }oder\textgreek[variant=ancient]{
ἐπεὶ εἴδομεν ἀλλήλους ὕδωρ ἀγένετο πολύ.}\par\end{continuousexamples}
\end{german}
\switchcolumn

\begin{japanese}
私たちが会って以来、雨が多かった：\textgreek[variant=ancient]{ἐξ οὗ }または
\textgreek[variant=ancient]{ἐπεὶ εἴδομεν ἀλλήλους ὕδωρ ἀγένετο πολύ.}
\end{japanese}
\switchcolumn*

\begin{german}
\item Wo sich statt \emph{sein} denken läßt \emph{gehen,} wird \textgreek[variant=ancient]{παρεῖναι
εἰς} angewandt.
\end{german}
\switchcolumn

\begin{japanese}
\emph{sein} の代わりに \emph{gehen} を想定できるところでは
\textgreek[variant=ancient]{παρεῖναι εἰς} を用いる。
\end{japanese}
\switchcolumn*

\begin{german}
\begin{continuousexamples}Sind Sie oft im Theater gewesen? \textgreek[variant=ancient]{ἦ
πολλάκις παρῆσθα εἰς τὸ θέατρον;}\par\end{continuousexamples}
\end{german}
\switchcolumn

\begin{japanese}
劇場へはよく行かれましたか？ \textgreek[variant=ancient]{ἦ πολλάκις παρῆσθα εἰς τὸ θέατρον;}
\end{japanese}
\switchcolumn*

\begin{german}
\item Indefinita werden nach Negationen gern negativ, \textgreek[variant=ancient]{πω}
jedoch bleibt unverändert.
\end{german}
\switchcolumn

\begin{japanese}
不定語は否定の後でしばしば否定形になるが、\textgreek[variant=ancient]{πω} は変化しない。
\end{japanese}
\switchcolumn*

\begin{german}
\item Ja = doch (franz. \textfrench{si!}) dem Unglauben oder mangelhaften
Glauben versichernd: \textgreek[variant=ancient]{ναί! }
\end{german}
\switchcolumn

\begin{japanese}
「はい（いや、そうだ）」＝不信や不十分な信頼に対して念押しする（仏 \textfrench{si!}）：
\textgreek[variant=ancient]{ναί!}
\end{japanese}
\switchcolumn*

\begin{german}
\item \emph{Zu, allzu} bleibt meist unübersetzt; z.\,B. Wir sind zu wenige
\textgreek[variant=ancient]{ὀλίγοι ἐσμέν,} du hast zu menig\footnote{\textjapanese{原文はmenig（誤記。wenigのこと）。}} geschrieben
\textgreek[variant=ancient]{ὀλίγον ἔγραψας.} \textgreek[variant=ancient]{Τὸ
ὕδωρ ψυχρὸν ὥστε λούσασθαί ἐστιν} (zu kalt). \textgreek[variant=ancient]{Νέοι
ἔτι ἐσμὲν ὥστε τοῦτ᾽ εἰδέναι} (zu jung, als daß wir wissen könnten).
\end{german}
\switchcolumn

\begin{japanese}
\emph{zu, allzu} は多くの場合訳さない。例えば「私たちは少なすぎる」
\textgreek[variant=ancient]{ὀλίγοι ἐσμέν,}「君は書く量が少なすぎる」
\textgreek[variant=ancient]{ὀλίγον ἔγραψας.}「水は冷たすぎて入浴できない」
\textgreek[variant=ancient]{Τὸ ὕδωρ ψυχρὸν ὥστε λούσασθαί ἐστιν}。
「私たちはまだ若すぎてそれを知り得ない」
\textgreek[variant=ancient]{Νέοι ἔτι ἐσμὲν ὥστε τοῦτ᾽ εἰδέναι}。
\end{japanese}
\switchcolumn*

\begin{german}
\emph{Nicht genug} \textgreek[variant=ancient]{ὀλίγος.} Er hat nicht
genug zu leben \textgreek[variant=ancient]{βίον ἔχει ὀλίγον.} Ich
habe nicht genug Geld \textgreek[variant=ancient]{ἀργύριον ἔχω ὀλίγον.}
\end{german}
\switchcolumn

\begin{japanese}
「十分でない」＝\textgreek[variant=ancient]{ὀλίγος.}
彼は生活するには足りない \textgreek[variant=ancient]{βίον ἔχει ὀλίγον.}
私は金が足りない \textgreek[variant=ancient]{ἀργύριον ἔχω ὀλίγον.}
\end{japanese}
\switchcolumn*

\begin{german}
\emph{Genug} = ausreichend wird adjectivisch meist durch \textgreek[variant=ancient]{ἱκανός}
ausgedrückt. Geld genug \textgreek[variant=ancient]{ἱκανὸν ἀργύριον.}
Ich denke, zwanzig Schüler sind genug \textgreek[variant=ancient]{ἱκανοὺς
νομίζω μαθηδὰς εἴκοσιν.}
\end{german}
\switchcolumn

\begin{japanese}
「十分な」（= ausreichend）は形容詞として主に \textgreek[variant=ancient]{ἱκανός} で表す。
十分な金 \textgreek[variant=ancient]{ἱκανὸν ἀργύριον.}
私は20人の生徒で十分だと思う \textgreek[variant=ancient]{ἱκανοὺς νομίζω μαθηδὰς εἴκοσιν.}
\end{japanese}
\switchcolumn*

\begin{german}
\emph{Genug} = in Menge \textgreek[variant=ancient]{οὐκ ὀλίγος.}
\end{german}
\switchcolumn

\begin{japanese}
「十分な量」＝\textgreek[variant=ancient]{οὐκ ὀλίγος.}
\end{japanese}
\switchcolumn*

\begin{german}
\item Ein anderer = noch ein weiterer \textgreek[variant=ancient]{ἕτερος;
}ein anderer = irgend welcher andere \textgreek[variant=ancient]{ἄλλος.}
\end{german}
\switchcolumn

\begin{japanese}
「もう一つの／さらに別の」＝\textgreek[variant=ancient]{ἕτερος;}「何か別の」＝\textgreek[variant=ancient]{ἄλλος.}
\end{japanese}
\switchcolumn*

\begin{german}
Ich war dort und viele andere \textgreek[variant=ancient]{ἐγὼ παρεγενόμην
καὶ ἕτεροι πολλοί.} Nun, es giebt ja andere gute Bücher genug \textgreek[variant=ancient]{ἀλλ᾽
ἔστιν ἔτερα νὴ Δία χρηστά βιβλία οὐκ ὀλίγα.}
\end{german}
\switchcolumn

\begin{japanese}
私はそこにいたし、ほかにも多くの人がいた \textgreek[variant=ancient]{ἐγὼ παρεγενόμην
καὶ ἕτεροι πολλοί.} まあ、ほかにも良い本は十分ある
\textgreek[variant=ancient]{ἀλλ᾽ ἔστιν ἔτερα νὴ Δία χρηστά βιβλία οὐκ ὀλίγα.}
\end{japanese}
\switchcolumn*

\begin{continuousitemline}
\begin{german}
Keine andere Sache \textgreek[variant=ancient]{οὐκ
ἄλλο πρᾶγμα.}
\end{german}
\switchcolumn

\begin{japanese}
ほかの事柄は何もない \textgreek[variant=ancient]{οὐκ ἄλλο πρᾶγμα.}
\end{japanese}
\switchcolumn*

\begin{german}
Wer sonst? \textgreek[variant=ancient]{τίς ἄλλος;}
\end{german}
\switchcolumn

\begin{japanese}
ほかに誰が？ \textgreek[variant=ancient]{τίς ἄλλος;}
\end{japanese}
\switchcolumn*
\end{continuousitemline}

\begin{german}
\item Immer noch = \textgreek[variant=ancient]{ἔτι καὶ νῦν,}
\end{german}
\switchcolumn

\begin{japanese}
なお今も＝\textgreek[variant=ancient]{ἔτι καὶ νῦν,}
\end{japanese}
\switchcolumn*

\begin{continuousitemline}
\begin{german}
noch welches \textgreek[variant=ancient]{ἄλλο,}
\end{german}
\switchcolumn

\begin{japanese}
さらにどれか＝\textgreek[variant=ancient]{ἄλλο,}
\end{japanese}
\switchcolumn*

\begin{german}
noch einige \textgreek[variant=ancient]{ἄλλοι,}
\end{german}
\switchcolumn

\begin{japanese}
さらに幾人か＝\textgreek[variant=ancient]{ἄλλοι,}
\end{japanese}
\switchcolumn*

\begin{german}
noch irgend einer \textgreek[variant=ancient]{ἄλλος τις.}
\end{german}
\switchcolumn

\begin{japanese}
さらに誰か＝\textgreek[variant=ancient]{ἄλλος τις.}
\end{japanese}
\switchcolumn*

\begin{german}
Hat er noch (sonstiges) Geld? \textgreek[variant=ancient]{ἆρ᾽ ἔχει
ἀργύριον ἄλλο;}
\end{german}
\switchcolumn

\begin{japanese}
彼はまだ（ほかに）金を持っているか？ \textgreek[variant=ancient]{ἆρ᾽ ἔχει
ἀργύριον ἄλλο;}
\end{japanese}
\switchcolumn*

\begin{german}
Er hat \emph{welches} \textgreek[variant=ancient]{ἔχει.}
\end{german}
\switchcolumn

\begin{japanese}
持っている \textgreek[variant=ancient]{ἔχει.}
\end{japanese}
\switchcolumn*
\end{continuousitemline}

\begin{german}
\item Ihr \emph{beiden} alten Herren \textgreek[variant=ancient]{ὦ \emph{δύο}
πρεσβύτα.}
\end{german}
\switchcolumn

\begin{japanese}
お二人の老紳士よ \textgreek[variant=ancient]{ὦ \emph{δύο} πρεσβύτα.}
\end{japanese}
\switchcolumn*

\begin{continuousitemline}
\begin{german}
\emph{Diese beiden} alten Herren hier \textgreek[variant=ancient]{τὼ
πρεσβύτα τώδε.}
\end{german}
\switchcolumn

\begin{japanese}
\emph{この二人の}老紳士 \textgreek[variant=ancient]{τὼ πρεσβύτα τώδε.}
\end{japanese}
\switchcolumn*

\begin{german}
\emph{Diese beiden} \textgreek[variant=ancient]{τώδε (ἄμφω).}
\end{german}
\switchcolumn

\begin{japanese}
\emph{この二人} \textgreek[variant=ancient]{τώδε (ἄμφω).}
\end{japanese}
\switchcolumn*
\end{continuousitemline}

\begin{german}
\textgreek[variant=ancient]{ἄμφω} verlangt stets den \emph{Dual}
des beigefügten Substantivs, \textgreek[variant=ancient]{ἀμφότερος}
steht meist mit seinem Substantiv im \emph{Plural.}
\end{german}
\switchcolumn

\begin{japanese}
\textgreek[variant=ancient]{ἄμφω} は常に、添えられる名詞に\emph{双数}を要求する。
\textgreek[variant=ancient]{ἀμφότερος} はたいてい名詞を\emph{複数}にして用いる。
\end{japanese}
\switchcolumn*

\begin{german}
\item allein (= allein für sich) \textgreek[variant=ancient]{αὐτός,}
\end{german}
\switchcolumn

\begin{japanese}
ひとりで（＝自分だけで）\textgreek[variant=ancient]{αὐτός,}
\end{japanese}
\switchcolumn*

\begin{continuousitemline}
\begin{german}
allein (= der einzige) \textgreek[variant=ancient]{μόνος.}
\end{german}
\switchcolumn

\begin{japanese}
唯一の＝\textgreek[variant=ancient]{μόνος.}
\end{japanese}
\switchcolumn*

\begin{german}
Wir sind allein (unter uns) \textgreek[variant=ancient]{αὐτοί ἐσμεν.}
\end{german}
\switchcolumn

\begin{japanese}
私たちは（内輪だけで）ひとりきりだ \textgreek[variant=ancient]{αὐτοί ἐσμεν.}
\end{japanese}
\switchcolumn*

\begin{german}
Wir sind die einzigen \textgreek[variant=ancient]{μόνοι ἐσμέν.}
\end{german}
\switchcolumn

\begin{japanese}
私たちだけが唯一の者だ \textgreek[variant=ancient]{μόνοι ἐσμέν.}
\end{japanese}
\switchcolumn*
\end{continuousitemline}

\begin{german}
Ich habe die (schriftliche) Arbeit allein gemacht \textgreek[variant=ancient]{αὐτὸς
ἐγὼ ταῦτα ἔγραψα.} Dagegen \textgreek[variant=ancient]{μόνος ἐγὼ ταῦτα
ἔγραψα} ich bin der Einzige, der diese Arbeit gemacht hat.
\end{german}
\switchcolumn

\begin{japanese}
私はその（筆記の）課題を自分だけでやった
\textgreek[variant=ancient]{αὐτὸς ἐγὼ ταῦτα ἔγραψα.}
これに対して \textgreek[variant=ancient]{μόνος ἐγὼ ταῦτα ἔγραψα} は、
この課題をやったのは私一人だけだ、という意味である。
\end{japanese}
\switchcolumn*

\begin{german}
\item Ich habe mehr \emph{von diesen} (z.\,B. Söhne) wie von jenen (Töchter)
\textgreek[variant=ancient]{πλείους ἔχω τούτους ἢ ἐκείνας} (doch auch
\textgreek[variant=ancient]{ἐκείνους ἢ ταύτας}).
\end{german}
\switchcolumn

\begin{japanese}
私は\emph{これら}（例えば息子）の方を、あれら（娘）より多く持っている
\textgreek[variant=ancient]{πλείους ἔχω τούτους ἢ ἐκείνας}
（ただし \textgreek[variant=ancient]{ἐκείνους ἢ ταύτας} とも言う）。
\end{japanese}
\switchcolumn*

\begin{german}
\item Wollen = Lust haben, sich entschließen \textgreek[variant=ancient]{ἐθέλειν.}
\end{german}
\switchcolumn

\begin{japanese}
「したい」＝気がある、決心する \textgreek[variant=ancient]{ἐθέλειν.}
\end{japanese}
\switchcolumn*

\begin{continuousitemline}
\begin{german}
Wollen = wünschen \textgreek[variant=ancient]{βούλεσθαι.}
\end{german}
\switchcolumn

\begin{japanese}
「望む」＝\textgreek[variant=ancient]{βούλεσθαι.}
\end{japanese}
\switchcolumn*

\begin{german}
Er hat keine Lust \textgreek[variant=ancient]{οὐκ ἐθέλει.}
\end{german}
\switchcolumn

\begin{japanese}
彼はその気がない \textgreek[variant=ancient]{οὐκ ἐθέλει.}
\end{japanese}
\switchcolumn*

\begin{german}
(Sehnlich) wünschen \textgreek[variant=ancient]{ἐπιθυμεῖν. }
\end{german}
\switchcolumn

\begin{japanese}
（切に）望む \textgreek[variant=ancient]{ἐπιθυμεῖν.}
\end{japanese}
\switchcolumn*

\begin{german}
Wollen = darüber sein \textgreek[variant=ancient]{μέλλειν.}
\end{german}
\switchcolumn

\begin{japanese}
「〜しようとしている」＝\textgreek[variant=ancient]{μέλλειν.}
\end{japanese}
\switchcolumn*

\end{continuousitemline}


\begin{german}
Wohin eilen sie? Ich will einen Brief zum Briefkasten tragen \textgreek[variant=ancient]{ποῖ
θεῖς; ἐπιστολὴν μέλλω φέρειν εἰς τὸ κιβώτιον (γραμματοκιβώτιον). Ich
will gehen εἶμι oder βαδιοῦμαι.} 
\end{german}
\switchcolumn

\begin{japanese}
どこへ急いでいるのですか？ 手紙を郵便箱へ持って行くところです
\textgreek[variant=ancient]{ποῖ θεῖς; ἐπιστολὴν μέλλω φέρειν εἰς τὸ κιβώτιον (γραμματοκιβώτιον).}
「私は行くつもりだ」は \textgreek[variant=ancient]{εἶμι} または \textgreek[variant=ancient]{βαδιοῦμαι.}
\end{japanese}
\switchcolumn*

\begin{german}
\begin{continuousitemline}Ich will gehen \textgreek[variant=ancient]{εἶμι}
oder \textgreek[variant=ancient]{βαδιοῦμαι.}\par\end{continuousitemline}
\end{german}
\switchcolumn

\begin{japanese}
私は行くつもりだ \textgreek[variant=ancient]{εἶμι} または \textgreek[variant=ancient]{βαδιοῦμαι.}
\end{japanese}
\switchcolumn*

\begin{german}
\item Wo ist dein Bruder? \textgreek[variant=ancient]{ποῦ ᾽σθ᾽ ὁ σὸς ἀδελφός;}
\end{german}
\switchcolumn

\begin{japanese}
君の兄弟はどこにいるのか？ \textgreek[variant=ancient]{ποῦ ᾽σθ᾽ ὁ σὸς ἀδελφός;}
\end{japanese}
\switchcolumn*

\begin{german}
\item Bei = franz. \textfrench{chez} \textgreek[variant=ancient]{παρά} mit
\textlatin{Dat.}
\end{german}
\switchcolumn

\begin{japanese}
「〜の家に／〜のところで」（仏 \textfrench{chez}）＝\textgreek[variant=ancient]{παρά}＋\textlatin{Dat.}
\end{japanese}
\switchcolumn*

\begin{german}
\begin{continuousitemline}Zu = franz. \textfrench{chez} \textgreek[variant=ancient]{παρά}
mit \textlatin{Acc.}\par\end{continuousitemline}
\end{german}
\switchcolumn

\begin{japanese}
「〜の家へ／〜のところへ」（仏 \textfrench{chez}）＝\textgreek[variant=ancient]{παρά}＋\textlatin{Acc.}
\end{japanese}
\switchcolumn*

\begin{german}
\item %
\begin{tabular}[t]{ccccc}
Mitnehmen, & ~ & mitbringen & ~ & (von Sachen) \textgreek[variant=ancient]{φέρειν,}\tabularnewline
,, &  & ,, &  & (von Personen) \textgreek[variant=ancient]{ἄγειν.}\tabularnewline
\end{tabular}
\end{german}
\switchcolumn

\begin{japanese}
\begin{tabular}[t]{ccccc}
持って行く、 & ~ & 持って来る & ~ & （物について）\textgreek[variant=ancient]{φέρειν,}\tabularnewline
,, &  & ,, &  & （人について）\textgreek[variant=ancient]{ἄγειν.}\tabularnewline
\end{tabular}
\end{japanese}
\switchcolumn*

\begin{continuousitemline}%
\begin{german}
Ich will das Buch mitbringen \textgreek[variant=ancient]{οἴσω
τὸ βιβλίον.}
\end{german}
\switchcolumn

\begin{japanese}
私はその本を持って来るつもりだ \textgreek[variant=ancient]{οἴσω τὸ βιβλίον.}
\end{japanese}
\switchcolumn*

\begin{german}
Ich will dich mit (zu ihm) nehmen \textgreek[variant=ancient]{ἄξω
σε παρ᾽ αυτόν.}
\end{german}
\switchcolumn

\begin{japanese}
私は君を（彼のところへ）連れて行くつもりだ \textgreek[variant=ancient]{ἄξω
σε παρ᾽ αυτόν.}
\end{japanese}
\switchcolumn*
\end{continuousitemline}

\begin{german}
\item Ich gehe (hin) \textgreek[variant=ancient]{βαδίζω,}
\end{german}
\switchcolumn

\begin{japanese}
私は（向こうへ）行く \textgreek[variant=ancient]{βαδίζω,}
\end{japanese}
\switchcolumn*

\begin{continuousitemline}
\begin{german}
ich komme (her) \textgreek[variant=ancient]{ἔρχομαι,}
\end{german}
\switchcolumn

\begin{japanese}
私は（こちらへ）来る \textgreek[variant=ancient]{ἔρχομαι,}
\end{japanese}
\switchcolumn*

\begin{german}
ich bin hergegangen \textgreek[variant=ancient]{ἐλήλυθα,}
\end{german}
\switchcolumn

\begin{japanese}
私はこちらへ来ている \textgreek[variant=ancient]{ἐλήλυθα,}
\end{japanese}
\switchcolumn*

\begin{german}
ich bin gekommen \textgreek[variant=ancient]{ἥκω,}
\end{german}
\switchcolumn

\begin{japanese}
私は来ている \textgreek[variant=ancient]{ἥκω,}
\end{japanese}
\switchcolumn*

\begin{german}
ich bin wieder da \textgreek[variant=ancient]{ἥκω,}
\end{german}
\switchcolumn

\begin{japanese}
私は戻って来ている \textgreek[variant=ancient]{ἥκω,}
\end{japanese}
\switchcolumn*

\begin{german}
bis ich wieder da bin \textgreek[variant=ancient]{μέχρι ἃν ἥκω,}
\end{german}
\switchcolumn

\begin{japanese}
私が戻って来るまで \textgreek[variant=ancient]{μέχρι ἃν ἥκω,}
\end{japanese}
\switchcolumn*

\begin{german}
ich gehe \emph{(weiter)} \textgreek[variant=ancient]{χωρῶ,}
\end{german}
\switchcolumn

\begin{japanese}
私は\emph{（先へ）進む} \textgreek[variant=ancient]{χωρῶ,}
\end{japanese}
\switchcolumn*

\begin{german}
ich will ihn \emph{besuchen} \textgreek[variant=ancient]{εἶμι (εἴσειμι)
ὡς αὐτόν,}
\end{german}
\switchcolumn

\begin{japanese}
私は彼を\emph{訪ねる}つもりだ \textgreek[variant=ancient]{εἶμι (εἴσειμι)
ὡς αὐτόν,}
\end{japanese}
\switchcolumn*

\begin{german}
ich werde kommen \textgreek[variant=ancient]{ἥξω.}
\end{german}
\switchcolumn

\begin{japanese}
私は来るだろう \textgreek[variant=ancient]{ἥξω.}
\end{japanese}
\switchcolumn*

\begin{german}
Ich will gehen, \emph{um} ihn zu befragen \textgreek[variant=ancient]{εἶμι
ἐρωτήσων αὐτόν.}
\end{german}
\switchcolumn

\begin{japanese}
私は彼に尋ねる\emph{ために}行くつもりだ \textgreek[variant=ancient]{εἶμι
ἐρωτήσων αὐτόν.}
\end{japanese}
\switchcolumn*

\begin{german}
Ich komme her, \emph{um} mit\textcompwordmark{}zuspeisen \textgreek[variant=ancient]{ἔρχομαι
δειπνήσων.}
\end{german}
\switchcolumn

\begin{japanese}
私は一緒に食事をする\emph{ために}こちらへ来る \textgreek[variant=ancient]{ἔρχομαι
δειπνήσων.}
\end{japanese}
\switchcolumn*

\begin{german}
\emph{aus}gehen \textgreek[variant=ancient]{θύραζε ἐξιέναι} oder
\textgreek[variant=ancient]{θ. βαδίζειν.}
\end{german}
\switchcolumn

\begin{japanese}
\emph{外へ}出る \textgreek[variant=ancient]{θύραζε ἐξιέναι} または
\textgreek[variant=ancient]{θ. βαδίζειν.}
\end{japanese}
\switchcolumn*
\end{continuousitemline}

\begin{german}
\item Die \emph{guten} Schüler \textgreek[variant=ancient]{οἱ ἀγαθοὶ τῶν
μαθητῶν.}
\end{german}
\switchcolumn

\begin{japanese}
\emph{優秀な}生徒たち \textgreek[variant=ancient]{οἱ ἀγαθοὶ τῶν
μαθητῶν.}
\end{japanese}
\switchcolumn*

\begin{german}
\begin{continuousitemline}Die guten \emph{Schüler} \textgreek[variant=ancient]{οἱ
ἀγαθοὶ μαθηταί.}\par\end{continuousitemline}
\end{german}
\switchcolumn

\begin{japanese}
優秀な\emph{生徒たち} \textgreek[variant=ancient]{οἱ ἀγαθοὶ μαθηταί.}
\end{japanese}
\switchcolumn*

\begin{german}
\item \emph{Da} kommt der junge Mann herbai! \textgreek[variant=ancient]{τὸ
μειράκιον \emph{τοδὶ} (τόδε) προσέρχεται!} 
\end{german}
\switchcolumn

\begin{japanese}
\emph{ほら}、若者がやって来る！ \textgreek[variant=ancient]{τὸ
μειράκιον \emph{τοδὶ} (τόδε) προσέρχεται!}
\end{japanese}
\switchcolumn*

\begin{german}
\item Ich habe \emph{nichts zu} essen \textgreek[variant=ancient]{οὐκ ἔχω
καταφαγεῖν.}
\end{german}
\switchcolumn

\begin{japanese}
私は\emph{食べるものを何も}持っていない \textgreek[variant=ancient]{οὐκ ἔχω
καταφαγεῖν.}
\end{japanese}
\switchcolumn*

\begin{german}
\item \emph{hier,} den Ort des \emph{Sprechenden} bezeichnend, haißt \textgreek[variant=ancient]{ἐνθάδε,}
\end{german}
\switchcolumn

\begin{japanese}
\emph{ここ}、すなわち\emph{話し手}のいる場所を指す場合は \textgreek[variant=ancient]{ἐνθάδε,}
\end{japanese}
\switchcolumn*

\begin{continuousitemline}
\begin{german}
hier (dem Ort des Sprechenden \emph{nahe})
\textgreek[variant=ancient]{ἐνταῦθα,}
\end{german}
\switchcolumn

\begin{japanese}
ここ（話し手の場所に\emph{近い}ところ）は \textgreek[variant=ancient]{ἐνταῦθα,}
\end{japanese}
\switchcolumn*

\begin{german}
hier (= \emph{an Ort und Stelle,} am Orte selbst) \textgreek[variant=ancient]{αὐτοῦ.}
\end{german}
\switchcolumn

\begin{japanese}
ここ（＝\emph{その場で}、まさにその場所で）は \textgreek[variant=ancient]{αὐτοῦ.}
\end{japanese}
\switchcolumn*
\end{continuousitemline}

\begin{german}
\item Jemanden kennen \textgreek[variant=ancient]{γιγνώσκειν τινά.}
\end{german}
\switchcolumn

\begin{japanese}
誰かを知っている \textgreek[variant=ancient]{γιγνώσκειν τινά.}
\end{japanese}
\switchcolumn*

\begin{german}
\item Zwar nicht groß, aber schön \textgreek[variant=ancient]{μέγας μὲν
οὔ, καλὸς δέ.}
\end{german}
\switchcolumn

\begin{japanese}
大きくはないが、美しい \textgreek[variant=ancient]{μέγας μὲν
οὔ, καλὸς δέ.}
\end{japanese}
\switchcolumn*

\begin{german}
\item Er hat eine breite Stirn \textgreek[variant=ancient]{πλατὺ ἔχει \emph{τὸ}
μέτωπον.}
\end{german}
\switchcolumn

\begin{japanese}
彼は広い額をしている \textgreek[variant=ancient]{πλατὺ ἔχει \emph{τὸ}
μέτωπον.}
\end{japanese}
\switchcolumn*

\begin{german}
Sie hat allerliebste Hände \textgreek[variant=ancient]{\emph{τὰς}
χεῖρας ἔχει παγκάλας.}
\end{german}
\switchcolumn

\begin{japanese}
彼女はとても愛らしい手をしている \textgreek[variant=ancient]{\emph{τὰς}
χεῖρας ἔχει παγκάλας.}
\end{japanese}
\switchcolumn*

\begin{german}
\item Beabsichtigen, gedenken \textgreek[variant=ancient]{ἐπινο\emph{εῖν}}
oder \textgreek[variant=ancient]{διανο\emph{εῖσθαι.}}
\end{german}
\switchcolumn

\begin{japanese}
意図する、考えている \textgreek[variant=ancient]{ἐπινο\emph{εῖν}}
または \textgreek[variant=ancient]{διανο\emph{εῖσθαι.}}
\end{japanese}
\switchcolumn*

\begin{german}
\item Ich lerne die Gedichte Homers \emph{auswendig} \textgreek[variant=ancient]{μανθάνω
τὰ Ὁμήρου ἔπη.}
\end{german}
\switchcolumn

\begin{japanese}
私はホメロスの詩を\emph{暗記して}学ぶ \textgreek[variant=ancient]{μανθάνω
τὰ Ὁμήρου ἔπη.}
\end{japanese}
\switchcolumn*

\begin{continuousitemline}
\begin{german}
Ich \emph{kann} die Ilias \emph{auswendig}
\textgreek[variant=ancient]{ἐπίσταμαι Ἰλιάδα.}
\end{german}
\switchcolumn

\begin{japanese}
私はイリアスを\emph{暗唱できる} \textgreek[variant=ancient]{ἐπίσταμαι Ἰλιάδα.}
\end{japanese}
\switchcolumn*

\begin{german}
Ich könnte die Odyssee \emph{auswendig hersagen} \textgreek[variant=ancient]{δυναίμην
ἂν Ὀδύσσειαν ἀπὸ στόματος εἰπεῖν. }
\end{german}
\switchcolumn

\begin{japanese}
私はオデュッセイアを\emph{暗唱して言う}ことができるだろう
\textgreek[variant=ancient]{δυναίμην ἂν Ὀδύσσειαν ἀπὸ στόματος εἰπεῖν.}
\end{japanese}
\switchcolumn*
\end{continuousitemline}

\begin{german}
\item Mein Vater hat mich gezwungen, die Odyssee auswendig zu lernen \textgreek[variant=ancient]{ὁ
πατὴρ ἠνάγκασέ με Ὀδύσσειαν \emph{μαθεῖν}} = that\textcompwordmark{}sächlich
mit dem Lernen zu Stande zu kommen; \textgreek[variant=ancient]{ἠνάγκασέ
με \emph{μανθάνειν}} bedeutet nur: er zwang mich, mit dem Lernen mich
zu beschäftigen, zu befassen, zu bemühen. 
\end{german}
\switchcolumn

\begin{japanese}
父は私にオデュッセイアを暗記するよう強いた
\textgreek[variant=ancient]{ὁ πατὴρ ἠνάγκασέ με Ὀδύσσειαν \emph{μαθεῖν}}
＝実際に学び終えるところまで行くこと。
\textgreek[variant=ancient]{ἠνάγκασέ με \emph{μανθάνειν}} は、
ただ「父は私に、学ぶことに取り組み、携わり、努めるよう強いた」という意味である。
\end{japanese}
\switchcolumn*

\begin{german}
\item \textgreek[variant=ancient]{Εὖ λέγει }er hat Recht.
\end{german}
\switchcolumn

\begin{japanese}
\textgreek[variant=ancient]{Εὖ λέγει} 彼の言うことは正しい。
\end{japanese}
\switchcolumn*

\begin{german}
\begin{continuousitemline}\textgreek[variant=ancient]{καλῶς λέγει}
er spricht gut.\par\end{continuousitemline}
\end{german}
\switchcolumn

\begin{japanese}
\textgreek[variant=ancient]{καλῶς λέγει} 彼は上手に話す。
\end{japanese}
\switchcolumn*

\begin{german}
\item Ich habe mehr Geld als du, aber Karl hat \emph{das} meiste \textgreek[variant=ancient]{ἐγὼ
μὲν ἀργύριον ἔχω πλέον ἢ σύ, πλεῖστον δὲ Κάρολος.}
\end{german}
\switchcolumn

\begin{japanese}
私は君より多くの金を持っているが、カールが\emph{最も}多く持っている
\textgreek[variant=ancient]{ἐγὼ μὲν ἀργύριον ἔχω πλέον ἢ σύ, πλεῖστον δὲ Κάρολος.}
\end{japanese}
\switchcolumn*

\begin{german}
\item Der Mann, \emph{dessen} Brief du liest \textgreek[variant=ancient]{ὁ
ἀνήρ, οὖ ἀναγιγνώσκεις \emph{τὴν} ἐπιστολήν.}
\end{german}
\switchcolumn

\begin{japanese}
君が読んでいる手紙の\emph{持ち主である}男
\textgreek[variant=ancient]{ὁ ἀνήρ, οὖ ἀναγιγνώσκεις \emph{τὴν} ἐπιστολήν.}
\end{japanese}
\switchcolumn*

\begin{german}
\begin{continuousitemline}Wessen Brief liest du? \textgreek[variant=ancient]{\emph{τὴν}
τίνος ἐπιστολὴν ἀναγιγνώσκεις;}\par\end{continuousitemline}
\end{german}
\switchcolumn

\begin{japanese}
君は誰の手紙を読んでいるのか？
\textgreek[variant=ancient]{\emph{τὴν} τίνος ἐπιστολὴν ἀναγιγνώσκεις;}
\end{japanese}
\switchcolumn*

\begin{german}
\item Setzest du deinen Hut auf? \textgreek[variant=ancient]{ἦ περιτίθεσαι
\emph{τὸν} πῖλον;}
\end{german}
\switchcolumn

\begin{japanese}
君は帽子をかぶるのか？ \textgreek[variant=ancient]{ἦ περιτίθεσαι
\emph{τὸν} πῖλον;}
\end{japanese}
\switchcolumn*

\begin{german}
\begin{continuousitemline}Zieh deine Stiefel aus! \textgreek[variant=ancient]{ἀποδύου
\emph{τὰς} ἐμβάδας!}\par\end{continuousitemline}
\end{german}
\switchcolumn

\begin{japanese}
長靴を脱ぎなさい！ \textgreek[variant=ancient]{ἀποδύου
\emph{τὰς} ἐμβάδας!}
\end{japanese}
\switchcolumn*

\begin{german}
Das Possessiv ist durch das Medium bereits ausgedrückt. 
\end{german}
\switchcolumn

\begin{japanese}
所有の意味は中動態によってすでに表されている。
\end{japanese}
\switchcolumn*

\begin{german}
\item Er wird dich von \emph{deinem} Augenleiden befreien \textgreek[variant=ancient]{ἀπαλλάξει
σε τῆς ὀφθαλμίας.}
\end{german}
\switchcolumn

\begin{japanese}
彼は君を\emph{君の}眼病から解放してくれるだろう
\textgreek[variant=ancient]{ἀπαλλάξει σε τῆς ὀφθαλμίας.}
\end{japanese}
\switchcolumn*

\begin{continuousitemline}
\begin{german}
Ein einziger Tag hat mir \emph{meinen}
ganzen Wohlstand geraubt \textgreek[variant=ancient]{μία ἡμέρα με
\emph{τὸν} πάντα ὄλβον ἀφείλετο.}
\end{german}
\switchcolumn

\begin{japanese}
たった一日が私から\emph{私の}全財産を奪った
\textgreek[variant=ancient]{μία ἡμέρα με \emph{τὸν} πάντα ὄλβον ἀφείλετο.}
\end{japanese}
\switchcolumn*

\begin{german}
Er hat mir \emph{mein} Geld gestahlen \textgreek[variant=ancient]{ὑπείλετό
μου τἀργύρια.}
\end{german}
\switchcolumn

\begin{japanese}
彼は私から\emph{私の}金を盗んだ \textgreek[variant=ancient]{ὑπείλετό
μου τἀργύρια.}
\end{japanese}
\switchcolumn*
\end{continuousitemline}

\begin{german}
Bei den Verben \emph{nehmen} und dergl. darf kein Possessiv übersetzt
werden, sobald die durch dasselbe bezeichnete Person bereits genannt
ist.
\end{german}
\switchcolumn

\begin{japanese}
\emph{取る}などの動詞では、所有によって示される人物がすでに言及されているなら、
所有代名詞を訳してはならない。
\end{japanese}
\switchcolumn*

\begin{german}
\item Brauchst du \emph{etwas?} \textgreek[variant=ancient]{δέει \emph{τίνος;}}
\end{german}
\switchcolumn

\begin{japanese}
\emph{何か}必要か？ \textgreek[variant=ancient]{δέει \emph{τίνος;}}
\end{japanese}
\switchcolumn*

\begin{german}
\begin{continuousitemline}Giebt es \emph{was} Neues? \textgreek[variant=ancient]{λέγεται
\emph{τί} καινόν;}\par\end{continuousitemline}
\end{german}
\switchcolumn

\begin{japanese}
\emph{何か}新しいことが言われているか？ \textgreek[variant=ancient]{λέγεται
\emph{τί} καινόν;}
\end{japanese}
\switchcolumn*

\begin{german}
\item Woher kommst du? \textgreek[variant=ancient]{πόθεν ἥκεις;} Aus dem
Garten \textgreek[variant=ancient]{ἐκ τοῦ κήπου.} Aus welchem? \textgreek[variant=ancient]{ἐκ
\emph{τοῦ} ποίου;}
\end{german}
\switchcolumn

\begin{japanese}
君はどこから来たのか？ \textgreek[variant=ancient]{πόθεν ἥκεις;}
庭から \textgreek[variant=ancient]{ἐκ τοῦ κήπου.} どの庭から？
\textgreek[variant=ancient]{ἐκ \emph{τοῦ} ποίου;}
\end{japanese}
\switchcolumn*

\begin{german}
Wenn \textgreek[variant=ancient]{ποῖος} auf einen mit Artikel versehenen
Gattungsnamen \textlatin{(Substantivum appellativum)} oder einen ihn
vertretenden Satz zurückweist, so nimmt es den Artikel an. Weg bleibt
der Artikel in der Regel nur dann, wenn \textgreek[variant=ancient]{ποῖος}
Prädicat ist. 
\end{german}
\switchcolumn

\begin{japanese}
\textgreek[variant=ancient]{ποῖος} が冠詞付きの普通名詞
\textlatin{(Substantivum appellativum)}、またはそれを代用する文を指し戻す場合、
冠詞を取る。冠詞が原則として落ちるのは、\textgreek[variant=ancient]{ποῖος} が述語である場合だけである。
\end{japanese}
\switchcolumn*

\begin{german}
\item \emph{Geld} in kleineren Summen \textgreek[variant=ancient]{ἀργύριον.}
\end{german}
\switchcolumn

\begin{japanese}
少額の\emph{金} \textgreek[variant=ancient]{ἀργύριον.}
\end{japanese}
\switchcolumn*

\begin{german}
\begin{continuousitemline}Geld = Kapitalien \textgreek[variant=ancient]{χρήματα.
}\par\end{continuousitemline}
\end{german}
\switchcolumn

\begin{japanese}
金＝資本 \textgreek[variant=ancient]{χρήματα.}
\end{japanese}
\switchcolumn*

\begin{german}
\item \textgreek[variant=ancient]{τάχα} ent\textcompwordmark{}spricht
genau dem in unserer Volks\textcompwordmark{}sprache üblichen \emph{am
Ende} (= schließlich, möglicher Weise)
\end{german}
\switchcolumn

\begin{japanese}
\textgreek[variant=ancient]{τάχα} は、私たちの民衆語でよく使う
\emph{am Ende}（＝結局、もしかすると）に正確に対応する。
\end{japanese}
\switchcolumn*

\begin{german}
\textgreek[variant=ancient]{ταχύ, ταχέως} schnell, bald,
\end{german}
\switchcolumn

\begin{japanese}
\textgreek[variant=ancient]{ταχύ, ταχέως} 速く、まもなく、
\end{japanese}
\switchcolumn*

\begin{german}

\textgreek[variant=ancient]{διὰ ταχέων} bald. 
\end{german}
\switchcolumn

\begin{japanese}
\textgreek[variant=ancient]{διὰ ταχέων} まもなく。
\end{japanese}
\switchcolumn*

\begin{german}
% 71
\item \emph{Unter} = zwischen drin \textgreek[variant=ancient]{ἐν,} z.\,B.
\textgreek[variant=ancient]{ἐν τοῖς Χριστιανοῖς πολλοί εἰσιν Ἰουδαῖοι.
ἐν νέοις ἀνὴρ γέρων.}
\end{german}
\switchcolumn

\begin{japanese}
\emph{〜の中に}＝その間に、の意味では \textgreek[variant=ancient]{ἐν,} 例えば
\textgreek[variant=ancient]{ἐν τοῖς Χριστιανοῖς πολλοί εἰσιν Ἰουδαῖοι.
ἐν νέοις ἀνὴρ γέρων.}
\end{japanese}
\switchcolumn*

\begin{german}
\item \emph{Nicht sonderlich} \textgreek[variant=ancient]{οὐ πάνυ.} Er strengt
sich nicht sonderlich an \textgreek[variant=ancient]{οὐ πάνυ σπουδάζει.}
\end{german}
\switchcolumn

\begin{japanese}
\emph{特に〜ではない} \textgreek[variant=ancient]{οὐ πάνυ.}
彼は特に熱心ではない \textgreek[variant=ancient]{οὐ πάνυ σπουδάζει.}
\end{japanese}
\switchcolumn*

\begin{german}
\item Die natürliche Stellung des Adverbs ist im Griechischen \emph{vor}
dem durch dasselbe zu bestimmenden Begriffe. Abweichung von dieser
Stellung dient zur Hervorhebung des Adverbs. Steht das Adverb mit
Nachdruck zuletzt, so ersetzt diese Stellung das deutsche \emph{und
zwar:} \textgreek[variant=ancient]{χάριν σωθέντες ὑπὸ σοῦ σοὶ ἂν ἔχοιμεν
δικαίως} (und zwar pflichtschuldigst). 
\end{german}
\switchcolumn

\begin{japanese}
ギリシア語で副詞の自然な位置は、それが限定する概念の\emph{前}である。
この位置から外れる場合は、副詞を強調するためである。
副詞が強勢を伴って最後に置かれると、この位置がドイツ語の
\emph{und zwar} に相当する：
\textgreek[variant=ancient]{χάριν σωθέντες ὑπὸ σοῦ σοὶ ἂν ἔχοιμεν δικαίως}
（しかも当然の義務として）。
\end{japanese}
\switchcolumn*

\begin{german}
\item Indirecte Ausrufesätze werden in der lateinischen Grammatik den indirecten
Fragesätzten gleichgestellt; im Griechischen unterscheiden sie sich
aber von den indirecten Fragesätzen dadurch, daß diese letzteren mit
dem indirecten oder directen Frageworte beginnen, die Ausrufesätze
hingegen mit dem Relativum, und zwar mit dem \emph{einfachen} Relativum.
\end{german}
\switchcolumn

\begin{japanese}
間接感嘆文は、ラテン文法では間接疑問文と同列に扱われる。
しかしギリシア語では間接疑問文と区別される。
すなわち、間接疑問文は間接または直接の疑問詞で始まるのに対し、
感嘆文は関係詞、それも\emph{単純な}関係詞で始まる。
\end{japanese}
\switchcolumn*

\begin{german}
\item Der Deutsche fragt: \emph{Wohin} setzt er sich? der Grieche: \emph{Wo?}
Wohin wollen wir uns setzen? \textgreek[variant=ancient]{ποῦ καθιζησόμεθα;}
\end{german}
\switchcolumn

\begin{japanese}
ドイツ人は「彼は\emph{\ruby{どこへ}{Wohin}}座るのか」と問うが、ギリシア人は「\emph{どこに}」と問う。
私たちはどこに座ろうか？ \textgreek[variant=ancient]{ποῦ καθιζησόμεθα;}
\end{japanese}
\switchcolumn*

\begin{german}
\item \emph{Alle Welt} \textfrench{(tout le monde)} heißt \textgreek[variant=ancient]{πάντες
ἄνθρωποι} (ohne Artikel).
\end{german}
\switchcolumn

\begin{japanese}
\emph{誰もかれも}（\textfrench{tout le monde}）は
\textgreek[variant=ancient]{πάντες ἄνθρωποι}（冠詞なし）という。
\end{japanese}
\switchcolumn*

\begin{german}
\item \emph{Um zu} wird gern durch \textgreek[variant=ancient]{βουλόμενος}
aus|gedrückt.
\end{german}
\switchcolumn

\begin{japanese}
\emph{〜するために} はしばしば \textgreek[variant=ancient]{βουλόμενος} で表される。
\end{japanese}
\switchcolumn*

\begin{german}
\item Ich habe bekommen = \textgreek[variant=ancient]{ἔχω,} z.\,B. ich habe
von meinem Vater 10 Mk. bekommen, \textgreek[variant=ancient]{δέκα
μάρκας ἔχω παρὰ τοῦ πατρός.}
\end{german}
\switchcolumn

\begin{japanese}
「受け取った」＝\textgreek[variant=ancient]{ἔχω,} 例えば、私は父から10マルクを受け取った、
\textgreek[variant=ancient]{δέκα μάρκας ἔχω παρὰ τοῦ πατρός.}
\end{japanese}
\switchcolumn*

\begin{german}
\item \emph{Lieber als} \ldots{} = eher als \ldots{} heißt \textgreek[variant=ancient]{μᾶλλον
ἢ }\ldots{}
\end{german}
\switchcolumn

\begin{japanese}
\emph{〜よりむしろ} \ldots{} ＝〜より先に \ldots{} は
\textgreek[variant=ancient]{μᾶλλον ἢ }\ldots{} という。
\end{japanese}
\switchcolumn*

\begin{german}
\item \emph{Vorhin} heißt \textgreek[variant=ancient]{τότε.} 
\end{german}
\switchcolumn

\begin{japanese}
\emph{先ほど} は \textgreek[variant=ancient]{τότε} という。
\end{japanese}
\switchcolumn*

\begin{german}
\item \textgreek[variant=ancient]{μέν} steht anderen Bindewörtern voran,
also nicht \textgreek[variant=ancient]{πολλοὶ γὰρ μὲν }\ldots{}\textgreek[variant=ancient]{,}
sondern \textgreek[variant=ancient]{πολλοὶ μὲν γὰρ }\ldots{}\textgreek[variant=ancient]{,}
ebenso \textgreek[variant=ancient]{μέν γε, μὲν δή }\ldots{}\textgreek[variant=ancient]{,
μὲν οὖν }\ldots{}\textgreek[variant=ancient]{, μέντοι.}
\end{german}
\switchcolumn

\begin{japanese}
\textgreek[variant=ancient]{μέν} は他の接続詞に先立つ。
従って \textgreek[variant=ancient]{πολλοὶ γὰρ μὲν }\ldots{}\textgreek[variant=ancient]{,}
ではなく \textgreek[variant=ancient]{πολλοὶ μὲν γὰρ }\ldots{}\textgreek[variant=ancient]{,}
同様に \textgreek[variant=ancient]{μέν γε, μὲν δή }\ldots{}\textgreek[variant=ancient]{,
μὲν οὖν }\ldots{}\textgreek[variant=ancient]{, μέντοι.}
\end{japanese}
\switchcolumn*

\begin{german}
\item Den bringlichen Imperativ, welchen wir durch \emph{so} (mach') \emph{doch}
ausdrücken, giebt der Grieche durch (das sehr oft und gern angewendete)
\textgreek[variant=ancient]{οὐ} mit Futurum, z.\,B. so schweig' doch!
\textgreek[variant=ancient]{οὐ σιγήσει;} Negation ist dabei \textgreek[variant=ancient]{μή,}
z.\,B. so mach' doch kein Gerede! \textgreek[variant=ancient]{οὐ μὴ
λαλήσεις; }so halte dich doch nicht auf! \textgreek[variant=ancient]{οὐ
μὴ διατρίψεις;}
\end{german}
\switchcolumn

\begin{japanese}
私たちが \emph{so} (mach') \emph{doch} で表す催促的命令法を、
ギリシア人は（非常によく、好んで用いられる）\textgreek[variant=ancient]{οὐ}＋未来で表す。
例えば「黙ってくれ！」\textgreek[variant=ancient]{οὐ σιγήσει;}
この場合の否定は \textgreek[variant=ancient]{μή} である。
例えば「そんなおしゃべりをしないでくれ！」\textgreek[variant=ancient]{οὐ μὴ λαλήσεις;}
「ぐずぐずしないでくれ！」\textgreek[variant=ancient]{οὐ μὴ διατρίψεις;}
\end{japanese}
\switchcolumn*

\begin{german}
\item Satzverbindungen wie folgende: \quotedblbase Wenn ich nach Dresden
komme und über die Brücke gehe, so sehe ich das Denkmal August des
Starken`` werden im Griechischen zerlegt in: \quotedblbase Wenn
ich nach Dresden komme, so sehe ich, wenn ich über die Brücke komme,
das Denkmal.`` Trotzdem gehen die \emph{beiden} Nebensätze dem Hauptsatze
voran. 
\end{german}
\switchcolumn

\begin{japanese}
次のような文結合、すなわち「私がドレスデンへ行き、橋を渡ると、アウグスト強王の記念碑が見える」は、
ギリシア語では「私がドレスデンへ行くと、橋のところへ来た時に、その記念碑が見える」と分解される。
それでも\emph{二つの}従属節は主節に先立つ。
\end{japanese}
\switchcolumn*

\begin{german}
\item Der gewöhnliche Ausdruck für \emph{\quotedblbase ich bitte``} ist
\textgreek[variant=ancient]{πρὸς (τῶν) θεῶν,} wofür auch \textgreek[variant=ancient]{πρὸς
τοῦ Διός} u. Ähnliches eintritt. \textgreek[variant=ancient]{πρὸς
θεῶν} ist keines\textcompwordmark{}wegs, wie gewöhnlich angegeben
wird, \emph{\quotedblbase Versicherung} bei den Göttern``, sondern
\emph{Bitt}formel.
\end{german}
\switchcolumn

\begin{japanese}
\emph{「お願いします」}の普通の表現は
\textgreek[variant=ancient]{πρὸς (τῶν) θεῶν,} であり、
\textgreek[variant=ancient]{πρὸς τοῦ Διός} などもその代わりに用いられる。
\textgreek[variant=ancient]{πρὸς θεῶν} は、ふつう説明されるような
\emph{「神々にかけての誓言」}では決してなく、\emph{懇願}の定型句である。
\end{japanese}
\switchcolumn*

\begin{german}
\item Es giebt nicht bloß, wie es nach den Grammatiken scheint, einen Irrealis
der Gegenwart und Irrealis der Vergangenheit (z.\,B. ich wäre (jetzt)
zufrieden, ich wäre (damals) zufrieden gewesen, wenn . . .), sondern
es muß auch einen Irrealis der \emph{Zukunft} geben. Ich sage z.\,B.:
\quotedblbase Wenn ich morgen in New-York wäre, würde ich mich an
dem Feste betheiligen,`` obgleich ich weiß, daß ich morgen unmöglich
dort sein kann. Diesen Irrealis der Zukunft drückt der Grieche im
Nebensatze durch \textgreek[variant=ancient]{εἰ} mit dem Optativ,
im regierenden Satze durch Optativ mit \textgreek[variant=ancient]{ἄν}
aus. 
\end{german}
\switchcolumn

\begin{japanese}
文法書からそう見えるように、現在の反実仮想と過去の反実仮想
（例えば「もし……なら、私は（今）満足しているだろう」「私は（その時）満足していただろう」）
だけがあるのではない。\emph{未来}の反実仮想も存在しなければならない。
例えば私は「もし明日ニューヨークにいたなら、その祭に参加するだろう」と言う。
明日そこにいることは不可能だと知っているにもかかわらずである。
この未来の反実仮想を、ギリシア人は従属節では
\textgreek[variant=ancient]{εἰ}＋希求法、主節では \textgreek[variant=ancient]{ἄν} を伴う希求法で表す。
\end{japanese}
\switchcolumn*
\end{enumerate}

\begin{german}
\emph{Anmerkung} In Beispielen, wie \textgreek[variant=ancient]{φαίη
δ᾽ ἂν ἡ θανοῦσα, εἰ φωνὴν λάβοι} steht also nicht der Optativ ungewöhnlich
für das Präteritum, sondern er bezeichnet regelrecht, wie in zahllosen
ähnlichen Fällen, den Irrealis der Zukunft: \quotedblbase wenn die
Verstorbene \emph{künftig einmal} wiederkäme, so würde sie es bestätigen.`` 
\end{german}
\switchcolumn

\begin{japanese}
\emph{注意} \textgreek[variant=ancient]{φαίη δ᾽ ἂν ἡ θανοῦσα, εἰ φωνὴν λάβοι}
のような例では、希求法が過去時制の代わりに異例に用いられているのではない。
むしろ無数の類例と同じく、規則どおり未来の反実仮想を示している。
すなわち「死者が\emph{将来いつか}戻って来るなら、彼女はそれを確証するだろう」という意味である。
\end{japanese}
\switchcolumn*

\end{paracol}

% 
% \begin{center}
% {\Huge{}\rule[1.2ex]{0.2\columnwidth}{1pt}}
% \par\end{center}{\Huge \par}
% 
% \clearpage{}



\leftskip=0.1in\parindent=-0.1in\tabcolsep=0pt
%\hangindent=2em\hangafter=2
%\titlespacing{\section}{0pt}{\parskip}{-\parskip}
\begin{paracol}{3}
	%\input{part-a.tex}
  %\setcounter{part}{1}
  %\setcounter{section}{20}
	%\input{part-b.tex}
	%\input{part-c.tex}
  \setcounter{part}{3}
  \setcounter{section}{39}
	\input{part-d.tex}
	\input{part-e.tex}
	\input{part-f.tex}
	\input{part-g.tex}
	%\input{part-h.tex}
	%\input{part-i.tex}
	% \input{lexicon.tex}
\end{paracol}

%\input{quotes.tex}

\cleartoevenpage{}
\pagestyle{empty}
\setlength{\parindent}{1\zw}
\setlength{\hangindent}{0pt}

%\section*{あとがき}

%本書を手に取って下さってありがとうございます。


\vspace*{\fill}

\begin{tcolorbox}[width=0.85\columnwidth,before=\par\smallskip\centering,after=\par,colback=white,coltext=black,]
  \begin{tabularx}{\columnwidth}{@{}r@{\hspace*{0.25\zw}：\hspace*{0.15\zw}}X@{}}
    \kintou{6\zw}{書名}&アッティカ方言で話しませんか? 日本語訳 第D部から第G部まで
    % \hfill{\includegraphics[trim={8cm 12cm 8cm 12cm},width=6em,valign=t]{./c100_logo_mono.png}}%
    \\
    \kintou{6\zw}{発行サークル}&ヒュアリニオス\newline\url{https://hyalinios.hatenadiary.com}%
    \\
    \kintou{6\zw}{発行責任者}&ざぎん\newline Twitter（現$\mathbb{X}$）: \texttt{@na4zagin3}%
    \\
    \kintou{6\zw}{発行日}&初版\quad \phantom{第1刷\quad }2026年8月16日
    \\
%    \kintou{6\zw}{印刷}&ちょ古っ都製本工房
%    \\
  \end{tabularx}
\end{tcolorbox}
\end{document}
